\documentclass{article}
\usepackage[utf8]{inputenc}
\usepackage{amsmath, amssymb}
\usepackage{graphicx}
\usepackage{hyperref}
\usepackage{geometry}
\geometry{a4paper, margin=1in}

\title{Analysis of Short-term Selection Experiments:}
\author{Unknown Authors}
\date{}

\begin{document}
\maketitle

\begin{abstract}

\end{abstract}

\section{1. Least-squares Approaches}
To consult the statistician after an experiment is finished is often merely to ask him to conduct a postmortem examination. He can perhaps say what the experiment died of. Fisher (1938)

This chapter examines the analysis of short-term selection experiments, whose duration is such that changes in allele frequency (from either selection or drift) are assumed to be relatively negligible. We restrict our attention here to the change in the mean under directional selection. The bulk of the chapter consists of developing least-squares (LS) estimates (and their corresponding standard errors) of realized heritabilities from the observed response. We also examine the empirical evidence for the goodness-of-fit of the breeder's equation to short-term artificial selection experiments. This is followed by a discussion of experimental and optimal designs for selection experiments, and we conclude with a detailed discussion of generation means analysis, the analysis of response when individuals from different cycles of selection are crossed (e.g., by using remnant seed).

Under the LS framework (LW Chapter 8), the only data required are the sample means (before and after selection) over generations, which are usually straightforward to obtain. Conversely, when one has access to the phenotypic values of essentially all of the individuals in the experiment, as well as their complete pedigree, more powerful mixed-model (MM) methods can be used. These are discussed in Chapter 19, while Chapter 20 considers the analysis of response in the much more complex setting of natural populations, where there is little control over either the nature of selection or the environment.

Finally, while our focus here is on phenotypic changes, we would be remiss if we did not point out that an area of very considerable interest is the genetic basis behind any such response. We previously discussed (Chapter 9; LW Chapter 15) a number of tools to aid in this dissection (e.g., QTL mapping using crosses between selected lines; tests of selection on specific loci/markers). The most recent extension of this quest are evolve & resequence (E&R) experiments (a term coined by Turner et al. [2011]), wherein a population is allowed to evolve under laboratory conditions and then genomic sequencing is applied in the search for target loci. See Long et al. (2015) for a recent review, and Kofler and Schlötterer (2013) and Bladwin-Brown et al. (2014) for design and power issues. Chapters 25 and 26 discuss the findings of such experiments.

The literature on selection experiments is truly massive, and our goal here is to introduce the important concepts rather than to exhaustively review all experiments. Reviews of selection experiments (mainly on metazoans) include Wilson (1977), Wright (1977), Robertson (1980), Mather (1983), Hill (1984), Sheridan (1988), Eisen (1989), Hill and Mackay (1989), Falconer (1992), Hill and Caballero (1992), Garland and Rose (2009), and Hill (2011). Fueled by the rapid advances in next-generation DNA sequencing, artificial selection experiments have undergone a renaissance, both in the form of experimental evolution studies (Kawecki et al. 2012) and in a growing number of experiments involving viral and microbial systems (Chapter 26). Our focus here is on goodness-of-fit to the breeder's equation, while in Chapters 25 and 26 we review empirical trends from long-term selection experiments.

\section{VARIANCE IN SHORT-TERM RESPONSE}
As we saw in Chapter 12, the means for a series of initially identical lines will inevitably diverge over time through the action of genetic drift, generating an among-line variance. The same is true for lines under selection, with this among-line variance being manifested

592

as variation in response (Figure 18.1). While the grand mean over a series of replicate lines under pure drift should remain essentially unchanged, the similar mean over a series of replicate selection lines should show a response over time, meaning that any particular realized response equals this expected response plus noise due to the evolutionary process of drift. Because artificial selection usually involves choosing a small number of parents to form the next generation, this among-line variance can be considerable. As we saw in Chapter 12, drift also changes the additive variance,  $ \sigma_A^2 $, within any given line. Our short-term assumption implies that  $ t/N_e \ll 1 $, in which case drift can generate a significant among-line variance in means but with little within-line change in  $ \sigma_A^2 $ (Chapters 11 and 12). Our focus on short-term response assumes that selection has only a minor role in changing the genetic variance over the course of the experiment. This is the standard infinitesimal model assumption used throughout Chapters 13–20 (and relaxed in Chapters 24–28). One complication that we initially ignore (addressing it later in the chapter) is that directional selection decreases the additive genetic variance by creating negative gametic-phase disequilibrium (Chapter 16).

\section{Expected Variance in Response Generated by Drift}
The expected variation among the means of replicate lines subjected to the same amount of selection was first considered by Prout (1962b) and then examined in detail by Hill (1971, 1972c, 1972d, 1974b, 1977a, 1980, 1986). The variation between the sample means,  $ \overline{z}_{t,i} $, of a series of replicate lines (in generation t) has two components: an evolutionary variance due to differences between the true means,  $ \mu_{t,i} $, generated by drift and selection, and a sampling (or residual) variance from estimating the true mean,  $ \mu_{t,i} $, by a sample mean,  $ \overline{z}_{t,i} $, based on  $ M_{t,i} $ individuals. Chapter 12 examined these components under pure drift, while here we consider the joint action of drift and selection.

With selection, the M individuals are not chosen at random with respect to their phenotypes, which decreases the among-line variance in population means relative to pure drift. This is partly countered by the fact that selection generates a lower effective population size than expected from drift alone (Chapters 3 and 26). Further, selection can increase the among-line variance relative to drift by changing allele frequencies more rapidly than expected from drift alone. It is these opposing (and potentially offsetting) changes in variance that led to suggestions that many of these effects will cancel out, leaving the pure-drift variance (with a suitably adjusted  $ N_{e} $) as an adequate approximation (Hill 1977a, 1980, 1986; Robertson 1977a; Nicholas 1980). Simulations (Robertson 1977a) and experimental results (Falconer 1973; López-Fanjul 1982) suggest that the pure-drift variance is a reasonable approximation.

Using the pure-drift results (Chapter 12), the sample mean in generation t for a particular realization (line) can be decomposed as

\begin{equation}
\overline{z}_{t}=\mu_{t}+e_{t}
\tag{18.1a}
\label{eq:18_1a}
\end{equation}

where  $ \mu_{t} $ is the true mean of this line and  $ e_{t} $ is the residual error in estimating this true mean from a population sample. The true mean can be further decomposed as

\begin{equation}
\mu_{t}=\mu+g_{t}+b_{t}
\tag{18.1b}
\label{eq:18_1b}
\end{equation}

where  $ g_t $ is the mean breeding value (relative to a base population value of  $ g_0 = 0 $) and  $ b_t $ is the mean environmental deviation in generation  $ t $, yielding an expected value for the sample mean of a random line in generation  $ t $ of

\begin{equation}
E(\overline{z}_{t})=\mu+E(g_{t})+b_{t}
\tag{18.2}
\label{eq:18_2}
\end{equation}

Under pure drift,  $ E(g_t) = 0 $, while under selection,  $ E(g_t) $ is given by the breeder's equation,  $ \sum_i^t h_i^2 S_i = th^2 S $, if the heritability and selection differential remain constant (given the usual

593

\begin{figure}[h]
\centering
\includegraphics[width=0.8\linewidth]{example-image}
\caption{Figure}
\end{figure}

\begin{figure}[h]
\centering
\includegraphics[width=0.8\linewidth]{example-image}
\caption{Figure}
\end{figure}

caveats summarized in Table 13.2). Assuming  $ b_t \sim (0, \sigma_b^2) $, the variance about the expected mean value is

\begin{equation}
\sigma_{\overline{{z}}}^{2}(t)=\sigma_{g}^{2}(t)+\sigma_{e}^{2}(t)+\sigma_{b}^{2}
\tag{18.3}
\label{eq:18_3}
\end{equation}

As discussed later in the chapter, comparison between contemporaneous populations, such as a selected and control line in the same generation or an upselected versus downselected line, can remove the shared environmental effect, and hence the  $ \sigma_{b}^{2} $ term.

594

Given our previous comments, we assume the evolutionary sampling variance,  $ \sigma_g^2(t) $, about the mean breeding value for selected lines is, to a good approximation, the same as that for lines under drift alone. If  $ M_0 $ individuals are initially sampled to form each line, then from Equation 12.1b

\begin{equation}
\sigma_{g}^{2}(t)=\left(\frac{1}{M_{0}}+2f_{t}\right)h^{2}\sigma_{z}^{2}
\tag{18.4}
\label{eq:18_4}
\end{equation}

The  $ M_{0} $ term accounts for variation in mean breeding value among lines in the founding generation, while  $ f_{t} $ (the amount of inbreeding at generation t) accounts for variation generated by subsequent drift. If population size remains constant

\begin{equation}
2f_{t}=2\left[1-\left(1-\frac{1}{2N_{e}}\right)^{t}\right]\simeq t/N_{e}\qquadfor t/N_{e}\ll1
\tag{18.5}
\label{eq:18_5}
\end{equation}

If different numbers of males  $ (N_{m}) $ and females  $ (N_{f}) $ are sampled or if N varies over time

\begin{equation}
2f_{t}\simeq\sum_{k=0}^{t-1}\left[\frac{1}{4N_{m}(k)}+\frac{1}{4N_{f}(k)}\right]
\tag{18.6}
\label{eq:18_6}
\end{equation}

The variance,  $ \sigma_e^2(t) $, in Equation 18.3, which is associated with estimating the true mean from a sample mean, depends on the relatedness among the  $ M_t $ individuals chosen to estimate the mean. If these individuals are unrelated, then  $ \sigma_e^2(t) = \sigma_z^2(t)/M_t $. If some selected individuals are related (such as from the same family), the positive covariance between them reduces the effective sample variance, with the exact form of  $ \sigma_e^2 $ depending on the distribution of family sizes within the sample. Hill (1971, 1980) shows that it is bounded by

\begin{equation}
\frac{\sigma_{z}^{2}-\sigma_{A}^{2}/2}{M_{t}}\leq\sigma_{e}^{2}(t)\leq\frac{\sigma_{z}^{2}}{M_{t}}
\tag{18.7}
\label{eq:18_7}
\end{equation}

which corresponds to the range where all individuals are from the same family (lower bound) to all being unrelated (upper bound). Conservatively choosing the upper bound, Equation 18.3 becomes

\begin{equation}
\sigma_{\overline{z}}^{2}(t)=\left(\frac{1}{M_{0}}+2f_{t}\right)h^{2}\sigma_{z}^{2}+\sigma_{b}^{2}+\sigma_{z}^{2}/M_{t}
\tag{18.8}
\label{eq:18_8}
\end{equation}

Because drift variance accumulates in each generation (via $f_t$ increasing each generation) while the other variance coefficients do not change, when $h^2\sigma_z^2 > \sigma_b^2$, the drift term in Equation 18.8 is expected to eventually dominate, usually after a few generations. Whereas Equation 18.8 describes the divergence between lines due to drift, drift also introduces a positive covariance between the means at different generations within a line (Equation 12.2).

\begin{equation}
\sigma(g_{t},g_{t^{\prime}})=\left(\frac{1}{M_{0}}+2f_{t}\right)h^{2}\sigma_{z}^{2}\quad for t<t^{\prime}
\tag{18.9}
\label{eq:18_9}
\end{equation}

Assuming that cross-generational environmental and residual effects are uncorrelated, so that  $ \sigma(b_t, b_{t'}) = \sigma(e_t, e_{t'}) = 0 $, then  $ \sigma(\overline{z}_t, \overline{z}_{t'}) = \sigma(g_t, g_{t'}) $.

The careful reader will notice that neither  $ h^2 $ or  $ \sigma_z^2 $ are indexed by generation, and hence they are assumed to be constant. This treatment is in keeping with our assumption that these parameters remain essentially unchanged over the short time course of the experiment. This is a major assumption and can be violated in several ways. First, changes in the underlying allele frequencies can change these variances (Chapters 5 and 24–28). If alleles of major effect are segregating, large changes in the variance can occur within a few generations. Further, when major alleles are segregating at low frequencies, the sampling to create lines from a base population can result in a substantial increase in the among-line variance over that

595

predicted by Equation 18.4, as these alleles are absent in some lines and overrepresented (relative to the base population) in others. This results in lines having a larger range of starting additive variances, thus increasing the variance in response (James 1970).

Second, directional selection generates negative gametic-phase disequilibrium, reducing the additive genetic variance within a line over the first few generations before approaching an equilibrium value (Chapter 16). We will deal with this shortly. Third, inbreeding due to finite population size reduces the additive genetic variance within lines. For a completely additive locus, the expected additive genetic variance (in the absence of mutational input) within a line in generation t is expressed by Equation 11.2

\begin{equation}
E\left[\sigma_{A}^{2}(t)\right]=\left(1-\frac{1}{2N_{e}}\right)^{t}\sigma_{A}^{2}(0)\simeq\left(1-\frac{t}{2N_{e}}\right)\sigma_{A}^{2}(0)\quad for\quad t\ll N_{e}
\end{equation}

where  $ \sigma_A^2(0) $ is the variance in the base population (the last expression follows from  $ (1-x)^t \approx 1 - xt $ for  $ |xt| \ll 1 $). If dominance or epistasis is present, the within-line variance can actually increase (for a time) under drift (Chapter 11), further inflating the among-line variance. Provided  $ t/2N_e \ll 1 $, the error introduced in Equations 18.8 and 18.9 by ignoring the reduction in the within-line variance due to inbreeding is small. Finally, when there is heritable variation in the environmental variance,  $ \sigma_E^2 $ is expected to increase under moderate to strong directional selection (Chapter 17), increasing  $ \sigma_z^2 $ and thus decreasing  $ h^2 $.

Despite all of these potential complications, those few experimental tests of the pure-drift approximation to artificial-selection data have found it to be fairly reasonable (Falconer 1973; López-Fanjul 1982). Mixed-model (REML/BLUP) methods developed in LW Chapters 26 and 28 (and reviewed here in Chapter 19), account for some of these concerns, but they require significantly more information (the phenotypic values of all individuals and their complete pedigree), which may be difficult to obtain.

\section{Variance in Predicted Response Versus Variance in Actual Response}
It is important not to confuse the previous expressions for variance in response (the variation about the mean response when the genetic parameters and selection differentials are known without error) with variance in the predicted response, which has (at least) two additional sources of variation. The first is the uncertainty produced by using estimates in place of the true values for the genetic and phenotypic variances. The second is uncertainty in the realized selection intensity. While selecting a preset fraction, p, of the population specifies the expected selection intensity, there is variation about this mean value for any particular realization (Chapter 14). For example, before selection is performed, we expect that truncation selection saving the uppermost 5% will have an average selection intensity of 2.06 (Example 14.1). However, when selecting the largest 5% from a small population (so that the number saved is small, e.g.,  $ N \leq 25 $ or so), its realized value can be significantly larger or smaller than the expectation, thus generating additional variation in the predicted response.

Several papers (Tai 1979; Knapp et al. 1989; Bridges et al. 1991) have presented confidence intervals for the expected selection response, but these simply focus on the variance generated by the uncertainty in the initial estimates of additive and phenotypic variances. They do not consider the additional evolutionary variance in response (discussed above), nor do they consider the variance in the particular realization of a prespecified selection intensity (although, retrospectively, we can deal with the latter by using the observed values).

\section{REALIZED HERITABILITIES}
The breeder’s equation,  $ R = h^2 S $, immediately suggests that heritability can be estimated as the ratio of the observed response to the observed selection differential

\begin{equation}
\widehat{h}_{r}^{2}=\frac{R}{S}
\tag{18.10}
\label{eq:18_10}
\end{equation}

596

Falconer (1954) referred to Equation 18.10 as the realized heritability, a term now more broadly defined to include any estimate of heritability based on the observed response to selection.

While one can use this approach to estimate  $ h^2 $, any complication in predicting response using the breeder's equation (Table 13.2) will usually make  $ \widehat{h}_r^2 $ a biased estimator. Turning this point around suggests that one test for the success of the breeder's equation is to compare how close realized heritabilities are to estimates based on resemblance between relatives in the unselected base population. If the breeder's equation generally provides an accurate model of selection response, we expect these two different estimates to be similar (i.e., within sampling error).

\begin{equation}
\widehat{h}_{r}^{2}=2R/S_{f}=2\cdot0.12/0.72=0.33
\end{equation}

This value was in good agreement with the estimated heritability based on mother-daughter regressions of  $ \widehat{h}^{2}=0.34\pm0.08 $.

\section{Estimators for Several Generations of Selection}
While the estimator given by Equation 18.10 is unambiguous for a single generation of selection, two different estimation approaches have been proposed for an experiment of a length of T > 1 generations. Both approaches are based on the cumulative selection response,  $ R_{C}(t) $, and cumulative selection differential,  $ S_{C}(t) $

\begin{equation}
S_{C}(t)=\sum_{i=1}^{t}S_{i}\qquad and\qquad R_{C}(t)=\sum_{i=1}^{t}R_{i}
\tag{18.11}
\label{eq:18_11}
\end{equation}

where  $ S_i = \overline{z}_i^* - \overline{z}_i $ and  $ R_i = \overline{z}_{i+1} - \overline{z}_i $ are the selection differential and single-generation response (respectively) for generation  $ i $, and  $ \overline{z}_i $ and  $ \overline{z}_i^* $ denoting the preselection and post-selection means in generation  $ i $.

The first approach, the simple ratio estimator of the realized heritability, is to use the total response divided by the total differential

\begin{equation}
\widehat{b}_{T}=\frac{R_{C}(T)}{S_{C}(T)}
\tag{18.12}
\label{eq:18_12}
\end{equation}

We use the notation throughout the chapter of  $ \hat{b}_x $ for a particular estimator, as the realized heritability is some multiple of this,  $ \hat{h}_r^2 = c \cdot \hat{b}_x $, where  $ c $ depends on the experimental design. For individual selection with only one sex selected,  $ c = 2 $ (Example 18.1), while  $ c = 1 $ when both sexes are selected.

597

The ratio estimator makes no assumption about a constant heritability, and instead returns an average value over the experiment. When  $ h^{2} $ remains constant, the breeder's equation predicts a linear cumulative response in R as a function of S. This observation leads to the second (and much more widely used) approach of estimating the realized heritability from the slope of the regression of cumulative response on cumulative selection differential

\begin{equation}
R_{C}(t)=b_{C}S_{C}(t)+e_{t}\quad for t=1,\cdots,T
\tag{18.13}
\label{eq:18_13}
\end{equation}

with  $ \hat{h}_{r}^{2}=b_{C} $ for individual selection on both sexes (Falconer 1954). Modifications of Equations 18.12 and 18.13 can be used for family selection (Chapter 21) and other designs, such as divergent selection (discussed at the end of this chapter). Because the expected response is zero if there is no selection, the regression line is constrained to pass through the origin and hence lacks an intercept term. Standard linear model approaches (LW Chapter 8) can be used to test for goodness-of-fit, such as testing whether a quadratic regression gives an improved fit (which suggests a changing heritability).

Recall from LW Chapter 8 that a regression estimator depends on the assumed covariance structure of the residuals. Most applications of the regression estimator (Equation 18.13) in the literature have assumed ordinary least-squares (OLS), which requires that the residuals,  $ e_{t} $, be homoscedastic and uncorrelated (LW Chapter 8). However, the careful reader will recall from Equations 18.4 and 18.9 that drift generates heteroscedastic and correlated residuals, thus requiring a generalized least-squares (GLS) solution (LW Chapter 8). We return to this issue after first examining the OLS estimator.

Recalling LW Equation 8.33a, the OLS estimator for the slope is  $ \mathbf{X}^T\mathbf{X}^{-1}\mathbf{X}^T\mathbf{y} $. Because the design matrix,  $ \mathbf{X} $, for the regression given by Equation 18.13 is simply the vector of cumulative selection differentials,  $ \mathbf{S} $, and  $ \mathbf{y} $ is the vector of cumulative responses,  $ \mathbf{R} $, it follows that the OLS estimate of the slope is given by

\begin{equation}
\widehat{b}_{C}(\mathrm{OLS})=\left(\mathbf{S}^{T}\mathbf{S}\right)^{-1}\mathbf{S}^{T}\mathbf{R}=\frac{\sum_{i=1}^{T}S_{C}(i)\cdot R_{C}(i)}{\sum_{i=1}^{T}S_{C}^{2}(i)}
\tag{18.14}
\label{eq:18_14}
\end{equation}

We will refer to this as the OLS regression estimator of the realized heritability. While this is the most common approach used in the literature, by assuming the wrong residual error structure, the OLS approach substantially underestimates the standard error for the slope. Richardson et al. (1968) and Irgang et al. (1985) noted that when the number of individuals,  $ M_t $, used to obtain the sample mean varies over generations, sample means based on more individuals contain more information, and thus have smaller residual errors. In this setting, the residual variances (roughly  $ \sigma_z^2/M_t $) are no longer homoscedastic, prompting these authors to suggest that a weighted least-squares method (e.g., LW Example 8.11) be used to correct for this. While it is an improvement over OLS, however, this approach still ignores the very significant impact from additional heteroscedasticity and correlations in the residual errors generated by drift.

The theory for estimating realized heritabilities (via a LS analysis) in populations with overlapping generations is less well-developed. As will be discussed in Volume 3, approaches assuming an asymptotic selection response are flawed in that many generations are required to reach a stable genetic structure when starting from an unselected base population. For nonasymptotic response, see Hill (1974a) and Johnson (1977a) for the relevant theory, and Atkins and Thompson (1986) for an example with Scottish Blackface sheep. Mixed-model approaches (Chapter 19) easily accommodate overlapping generations, but they require the full pedigree of all measured individuals used in making selection decisions. This is often quite feasible with selection on large domesticated animals, and it is becoming increasingly feasible in more general settings using dense-marker estimates of relationships (Chapter 20).

598

\begin{figure}[h]
\centering
\includegraphics[width=0.8\linewidth]{example-image}
\caption{Figure}
\end{figure}

\begin{table}[h]
\centering
\begin{tabular}{|c|c|}
\hline
Cell 1 & Cell 2 \\
\hline
Cell 3 & Cell 4 \\
\hline
\end{tabular}
\caption{Table Placeholder}
\end{table}

The ratio estimate of the realized heritability (Equation 18.12) is  $ \hat{h}_r^2 = 2.63/12.51 = 0.2102 $. The regression, forced through the origin, of cumulative response ( $ R_C $) on cumulative selection differential ( $ S_C $) is plotted in Figure 18.2. Equation 18.14 yields an OLS regression estimator of the realized heritability of

\begin{equation}
\widehat{h}_{r}^{2}=\widehat{b}_{C}(O L S)=\frac{\sum_{i=1}^{5}S_{C}(i)\cdot R_{C}(i)}{\sum_{i=1}^{5}S_{C}^{2}(i)}=\frac{78.96}{350.45}=0.2245
\end{equation}

\section{Weighted Least-squares Estimates of Realized Heritability}
Ordinary least-squares regression assumes that the residuals are homoscedastic and uncorrelated, yielding a covariance matrix for the vector of residuals, e, as  $ \operatorname{Var}(e) = \sigma_e^2 \cdot \mathbf{I} $ (LW Chapter 8). Genetic drift causes the residual covariance structure of the regression given by Equation 18.13 to depart significantly from this simple form. In particular, the residual variance increases with time (Equation 18.8) and residuals from different generations within a given line are correlated (Equation 18.9). Thus,  $ \operatorname{Var}(e) $ is not simply  $ \sigma_e^2 \cdot \mathbf{I} $, and generalized least-squares (GLS) must be used to account for the overall covariance structure. From LW Equation 8.34, the GLS estimator of the regression slope is given by

\begin{equation}
\widehat{b}_{C}(\mathrm{GLS})=(\mathbf{S}^{T}\mathbf{V}^{-1}\mathbf{S})^{-1}\mathbf{S}^{T}\mathbf{V}^{-1}\mathbf{R}
\tag{18.15a}
\label{eq:18_15a}
\end{equation}

599

where V is the variance-covariance matrix associated with selection response, with elements

\begin{equation}
V_{i j}=\sigma_{e}(i,j)=\sigma\left[R_{C}(i),R_{C}(j)\right]
\tag{18.15b}
\label{eq:18_15b}
\end{equation}

The elements of V can be obtained from the pure-drift approximation, with the variances,  $ V_{ii} $, given by Equation 18.8 and the covariances,  $ V_{ij} $, by Equation 18.9, yielding

\begin{equation}
V_{ij}=\left\{\begin{array}{ll}\left(\frac{1}{M_{0}}+2f_{i}\right)h^{2}\sigma_{z}^{2}+\sigma_{z}^{2}/M_{i}&for i=j\\\left(\frac{1}{M_{0}}+2f_{i}\right)h^{2}\sigma_{z}^{2}&for i<j\end{array}\right.
\tag{18.15c}
\label{eq:18_15c}
\end{equation}

where  $ 2f_i \simeq i/N_e $ (Equation 18.5). If differences in environmental values across generations are not accommodated for by the design, then  $ V_{ii} $ has an additional term,  $ \sigma_b^2 $, reflecting random changes in the environment that influence the mean trait value (as discussed below, the use of control lines can often remove this term). Even though OLS assumes an incorrect residual structure, it still provides an unbiased estimate of  $ b_C $. However, OLS significantly underestimates the standard error of  $ b_C $, and it is for this reason that GLS estimators should be used whenever possible (Hill 1971, 1972c, 1972d, 1974b, 1977a, 1980, 1986).

Example 18.3. Computing the GLS regression using the data from Example 18.2 requires the variance-covariance matrix,  $ \mathbf{V} $, of the residuals, which is obtained using the pure-drift approximation. From Example 18.2,  $ M = 100 $ and  $ N = 20 $, while the estimated phenotypic variance is  $ \text{Var}(z) = 3.293 $ (Mackay, personal communication). Assuming  $ N_e = N $, Equation 18.5 yields  $ f_i \simeq i/20 $. Further assuming that both initial sampling and among-generation environmental effects can be ignored (i.e.,  $ M_0 \gg 1 $ and  $ \sigma_b^2 \simeq 0 $), Equation 18.8 gives the variance associated with the response in generation  $ i $ as

\begin{equation}
V_{ii}=2f_{i}h^{2}\sigma_{z}^{2}+\sigma_{z}^{2}/M_{i}=\left(\frac{i}{20}\right)h^{2}\cdot3.293+\frac{3.293}{100}=0.1647\cdot(i\cdot h^{2}+0.2)
\end{equation}

Similarly, Equations 18.9 and 18.5 give the covariance between generations as

\begin{equation}
V_{ij}=\left(\frac{i}{N}\right)h^{2}\sigma_{z}^{2}=i\cdot h^{2}\cdot0.1647\qquadfor i<j
\end{equation}

The resulting covariance matrix becomes

\begin{equation}
\mathbf{V}=0.1647\cdot\begin{pmatrix}h^{2}+0.2&h^{2}&h^{2}&h^{2}&h^{2}\\h^{2}&2h^{2}+0.2&2h^{2}&2h^{2}&2h^{2}\\h^{2}&2h^{2}&3h^{2}+0.2&3h^{2}&3h^{2}\\h^{2}&2h^{2}&3h^{2}&4h^{2}+0.2&4h^{2}\\h^{2}&2h^{2}&3h^{2}&4h^{2}&5h^{2}+0.2\end{pmatrix}
\end{equation}

Because  $ \mathbf{V} $ is a function of the unknown heritability, estimation is an iterative process. Starting with some initial estimate of  $ h^2 $, each new  $ h_r^2 $ estimate is used to update  $ \mathbf{V} $ in subsequent iterations until convergence occurs. Using the ratio estimator  $ h^2 = 0.21 $ (Equation 8.2) as the starting value, Equation 18.15 gives a first estimate as

\begin{equation}
\widehat{b}_{C}(GLS)^{(1)}=(\mathbf{S}^{T}\mathbf{V}^{-1}\mathbf{S})^{-1}\mathbf{S}^{T}\mathbf{V}^{-1}\mathbf{R}=0.2222
\end{equation}

Substituting this new estimate of  $ h^2 $ into  $ \mathbf{V} $ gives, upon a second iteration,  $ \widehat{b}_C(GLS)^{(2)} = 0.2221 $, which remains unchanged in subsequent iterations.

600

\section{Standard Errors of Realized Heritability Estimates}
The final pieces of statistical machinery necessary for assessing the success of the breeder's equation are the standard errors associated with the different realized heritability estimators. Consider first the realized heritability estimated from the unweighted (OLS) regression (Equation 18.14). Recalling LW Equation 8.33b for the variance for an OLS estimator

\begin{equation}
\mathrm{Var}\left[\widehat{b}_{C}(\mathrm{OLS})\right]=\sigma_{e}^{2}\left(\mathbf{X}^{T}\mathbf{X}\right)^{-1}=\sigma_{e}^{2}\left(\mathbf{S}^{T}\mathbf{S}\right)^{-1}=\sigma_{e}^{2}\bigg/\sum_{i=1}^{T}S_{C}^{2}(i)
\tag{18.16a}
\label{eq:18_16a}
\end{equation}

The residual variance,  $ \sigma_{e}^{2} $, can be estimated from the residual sums of squares divided by the degrees of freedom (see LW Chapter 8). For an experiment lasting for T generations

\begin{equation}
\widehat{\sigma}_{e}^{2}=\frac{1}{T-1}\sum_{i=1}^{T}\widehat{e}_{i}^{2}=\frac{1}{T-1}\sum_{i=1}^{T}\left(R_{C}(i)-\widehat{h}_{r}^{2}S_{C}(i)\right)^{2}
\tag{18.16b}
\label{eq:18_16b}
\end{equation}

As mentioned, because the OLS estimator assumes that residuals are uncorrelated and have equal variances (both of which are incorrect under drift), it significantly underestimates the correct variance (Example 18.4). The GLS regression estimator (Equation 18.15) avoids these problems by properly accounting for the residual variance-covariance structure generated by drift. From standard GLS theory (LW Equation 8.35)

\begin{equation}
\mathrm{Var}\left[\widehat{b}_{C}(\mathrm{GLS})\right]=(\mathbf{S}^{T}\mathbf{V}^{-1}\mathbf{S})^{-1}
\tag{18.17}
\label{eq:18_17}
\end{equation}

As was done above, the pure-drift approximation is used to obtain the elements of V, with  $ \widehat{h}_{r}^{2} $ used in place of  $ h^{2} $.

Finally, consider the variance for the ratio estimator  $ b_T $, the total response to total selection (Equation 18.12). Because  $ \text{Var}(y/c) = \text{Var}(y)/c^2 $ for a constant, c, we have

\begin{equation}
\mathrm{Var}(\widehat{b}_{T})=\frac{\mathrm{Var}\left[R_{C}(T)\right]}{S_{C}^{2}(T)}\simeq\frac{(T/N)\widehat{h}_{r}^{2}\sigma_{z}^{2}+\sigma_{z}^{2}/M_{T}}{S_{C}^{2}(T)}
\tag{18.18}
\label{eq:18_18}
\end{equation}

The numerator (the variance in response in generation $T$) follows from the pure-drift approximation (Equation 18.8), assuming that initial sampling can be ignored (e.g., $M_0 \gg 1$) and there is no significant between-generation environmental variance ($\sigma_b^2 = 0$). Hill (1972c, 1972d) noted that $\widehat{b}_C(\mathrm{GLS})$ is generally a slightly better estimator (i.e., returning a smaller standard error) than $\widehat{b}_T$ when $h^2$ is small, while $\widehat{b}_T$ is a slightly better estimator when $h^2$ or the number of generations is large.

Example 18.4. Using the data from Examples 18.2 and 18.3, let us compare the standard errors associated with the three different realized heritability estimates. Consider the unweighted regression estimator,  $ \widehat{b}_{C}(\mathrm{OLS}) $, first. The residual sum of squares is

\begin{equation}
\sum_{i=1}^{T}\left(R_{C}(i)-\widehat{h}_{r}^{2}S_{C}(i)\right)^{2}=0.091
\end{equation}

yielding an estimated residual variance of  $ \widehat{\sigma}_{e}^{2}=0.091/4=0.0228 $. Equation 18.16a then yields

\begin{equation}
\mathrm{Var}\left[\widehat{b}_{C}(\mathrm{OLS})\right]=\sigma_{e}^{2}\bigg/\sum_{i=1}^{T}S_{C}^{2}(i)=\frac{0.0228}{350.45}=0.0000649
\end{equation}

601

Taking the square root yields a standard error of 0.0081.

Turning to the estimate,  $ b_{T} $, based on the total response to total selection, Equation 18.18 returns

\begin{equation}
\mathrm{Var}(\widehat{b}_{T})=\frac{(5/20)\cdot0.21\cdot3.292+0.03292}{12.51^{2}}=0.00132
\end{equation}

for a standard error of 0.0363. Finally, if we substitute the GLS estimate of  $ \widehat{h}_r^2 = 0.222135 $ (Example 18.3) into  $ \mathbf{V} $, Equation 18.17 will yield a variance for this estimate of  $ (\mathbf{S}^T \mathbf{V}^{-1} \mathbf{S})^{-1} = 1/790.4 $, for a standard error of 0.0356.

To summarize, the three approaches give extremely similar estimates.

\begin{equation}
\begin{array}{r l r l}&{\mathrm{U n w e i g h t e d~l e a s t-s q u a r e s~r e g r e s s i o n},\widehat{b}_{C}(\mathrm{O L S})}&{}&{\widehat{h}_{r}^{2}=0.2245\pm0.0081}\end{array}
\end{equation}

\begin{equation}
\begin{array}{r l r l r l r}&{}&{\mathrm{T o t a l~r e s p o n s e/t o t a l~d i f f e r e n t i a l},\widehat{b}_{T}}&&{\quad}&{\widehat{h}_{r}^{2}=0.2102\pm0.0363}\end{array}
\end{equation}

\begin{equation}
\begin{array}{r l r l r l}&{}&{\mathrm{W e i g h t e d~l e a s t-s q u a r e s~r e g r e s s i o n},\widehat{b}_{C}(\mathrm{G L S})}&{}&{\widehat{h}_{r}^{2}=0.2221\pm0.0356}\end{array}
\end{equation}

Note that the difference in the standard error between  $ \widehat{b}_{C} $(GLS) and  $ \widehat{b}_{T} $ is very small, with  $ \widehat{b}_{T} $ considerably more straightforward to compute. Also note that the unweighted regression greatly underestimates the standard error, making it four-fold smaller than the other two standard errors.

\section{Power: Estimation of  $ h^{2} $ from Relatives or Selection Response?}
Estimates of heritability from relatives are based on variance components (e.g., LW Chapters 17 and 18), while estimates from selection response are based on ratios of means. Thus, one might expect that realized heritability estimates may have greater power than those based on relatives. Is this indeed the case? Consider the sampling variance for the simple ratio estimator (Equation 18.18)

\begin{equation}
\mathrm{Var}(\widehat{h}_{r}^{2})=\frac{\mathrm{Var}\left[R_{C}(T)\right]}{S_{C}^{2}(T)}\simeq\frac{(T/N)\widehat{h}_{r}^{2}\sigma_{z}^{2}+\sigma_{z}^{2}/M}{S_{C}^{2}(T)}=\frac{(T/N)\widehat{h}_{r}^{2}+1/M}{S_{C}^{2}(T)/\sigma_{z}^{2}}
\end{equation}

where T is the number of generations and M individuals are measured per generation, with the most extreme N of these allowed to reproduce. Suppose that the strength of selection is the same in each generation  $ (S_i = S) $, in which case we have that

\begin{equation}
S_{C}^{2}(T)/\sigma_{z}^{2}=T^{2}\;S^{2}/\sigma_{z}^{2}=T^{2}\cdot\bar{\imath}^{2}
\end{equation}

and hence

\begin{equation}
\mathrm{Var}(\widehat{h}_{r}^{2})\simeq\frac{(T/N)\widehat{h}_{r}^{2}+1/M}{T^{2}\cdot\overline{\imath}^{2}}=\frac{(1/TN)\widehat{h}_{r}^{2}+1/(T^{2}M)}{\overline{\imath}^{2}}\simeq\frac{\widehat{h}_{r}^{2}}{TN\overline{\imath}^{2}}
\tag{18.19}
\label{eq:18_19}
\end{equation}

where the last step follows by ignoring the  $ 1/(T^{2}M) $ term.

Turning to relative-based estimators of  $ h^{2} $, the optimal midparent-offspring design ( $ N_{mp} $ sets of parents each with a single offspring) has (LW Equation 17.11b) a standard error of

\begin{equation}
\mathrm{SE}(h_{mp}^{2})=\left(\frac{2-h^{4}}{N_{mp}}\right)^{1/2}
\end{equation}

This gives a ratio of sampling variances for the two estimators of

\begin{equation}
\frac{\mathrm{Var}(h_{mp}^{2})}{\mathrm{Var}(\widehat{h}_{r}^{2})}\simeq\frac{(2-h^{4})/N_{mp}}{h^{2}/(TN\bar{\imath}^{2})}=\left(\frac{TN}{N_{mp}}\right)\left[\bar{\imath}^{2}\left(\frac{2-h^{4}}{h^{2}}\right)\right]
\tag{18.20}
\label{eq:18_20}
\end{equation}

In applying Equation 18.20 to determine which method ( $ N_{mp} $ midparent-offspring pairs versus TN total adults over the course of a selection experiment) is more accurate, a few

602

points are in order. In a selection experiment, one measures  $ M $ adults in each generation, where  $ N $ of these reproduce. Hence, if the limiting factor for an investigator is in raising offspring, then the comparison of  $ N_{mp} $ midparents with  $ TN $ adults that are allowed to reproduce is a fair one. However, if the limiting factor is in the actual measurement of individuals, then the fair comparison should be designs based on  $ TM $ measured adults versus  $ 3N_{mp} $ midparents (as the latter involves two adults and one offspring measure per pair). As Example 18.5 illustrates, even when  $ T $ and  $ M $ are fixed, the choice of  $ N $ is important, as the standard error (Equation 18.19) is minimized when the product,  $ N\ \bar{v}^2 $, is maximized. There is a tradeoff in this maximization, as increasing selection decreases  $ N $ but increases  $ \bar{v} $ (which is a nonlinear function of  $ N/M $).

There are a few final caveats in using a realized heritability estimator in place of a more traditional relative-based one. Clearly, if any of the assumptions of the breeder's equation (Table 13.2) are violated, we obtain a potentially biased estimate. A more subtle issue is the target population for the heritability estimate—is this the entire population or some subset? An example of why this distinction may matter derives from the work of Skibinski and Shereif (1989). These authors initiated lines for selection on sternopleural bristle number in Drosophila melanogaster by taking parents from the central part of the distribution versus taking parents with extremely high or extremely low bristle number. The lines founded from the central part of the distribution had higher heritabilities, and larger responses, than lines founded from the extremes of the distribution. This is not expected under the infinitesimal model (Chapter 24), but it is expected if there are alleles of modest (or greater) effect segregating in the population. Thus, the subtle point is that realized heritability estimates may actually be sampling a different part of the population than relative-based estimators (this was also stressed in Chapter 6; see Equations 6.39 and 6.40). For example, if we sample rare alleles of large effect, these will rapidly increase in selected populations, thus inflating the heritability.

Example 18.5. Consider a trait with a true heritability of  $ h^2 = 0.30 $. This trait is expensive to measure, while offspring are inexpensive to rear. With the resources to measure 3000 individuals, two possible designs are to measure 1000 midparent-offspring trios or to measure a total of 3000 individuals over the course of a selection experiment. The expected standard error for the midparent-offspring regression (LW Equation 17.11b) is

\begin{equation}
\mathrm{SE}(h_{mp}^{2})=\sqrt{\frac{2-h^{4}}{N_{mp}}}=\sqrt{\frac{2-0.30^{2}}{1000}}=0.044
\end{equation}

For a T-generation selection experiment, we have  $ T \cdot M = 3000 $ as the design constraint. Thus, if we measure M initial parents to select the most extreme N, and do this for six generations (the parents = generation 0, followed by selected generations 1, 2, 3, 4, 5), with  $ 3000/6 = 500 $ individuals being measured in each generation.

Consider three different designs, with N = 250 (50% selected), 100 (20%), or 50 (10%). Applying Equation 14.3a, these correspond to selection intensities (not correcting for finite population size) of  $ \bar{\imath} = 0.80 $, 1.40, and 1.75, respectively. From Equation 18.19, the standard errors for  $ h_{r}^{2} $ are given by

\begin{equation}
\sqrt{\frac{h^{2}}{TN\bar{\tau}^{2}}}=\sqrt{\frac{0.3}{5N\bar{\tau}^{2}}}
\end{equation}

or 0.0194, 0.0175, and 0.0197 (for N = 250, 100, and 50, respectively). Thus, while all three selection-experiment designs have a smaller standard error than the midparent-offspring estimate, the design with N = 100 is the best among the three candidates. Trial and error yields an optimal N of 135, resulting in a standard error of 0.0172. Similar calculations can also be performed assuming experiments of different length (with T varying).

603

\begin{figure}[h]
\centering
\includegraphics[width=0.8\linewidth]{example-image}
\caption{Figure}
\end{figure}

\section{Empirical Versus Predicted Standard Errors}
While the standard error in response can be directly estimated from the among-line variance over a series of replicate lines subjected to identical selection, this approach is generally not cost-effective given the large standard error on such estimates (Chapter 12). As a result, the standard errors reported for most experiments use the pure-drift approximation discussed previously. How closely do the empirical (observed) standard errors match those predicted by the drift approximation? Unfortunately, very few experiments have addressed this issue, and those few that do often use OLS-generated standard errors (which are too small) to compare the fit with the observed among-line variance.

Perhaps the most extensive study is that of Falconer (1973), who performed two-way selection for 6-week weight in mice (Figure 18.3). Six replicate sets of lines were used. Each set consisted of a line selected for larger size, a line selected for smaller size, and an unselected control, for a total of 18 lines in the entire experiment. Realized heritabilities were estimated (by OLS regression) using the three different contrasts available within each replicate set: large versus control (High), small versus control (Low), and large versus small (Divergent).

In addition to these three estimates of  $ \hbar_{r}^{2} $ for each of the six replicates, Falconer also considered two combined estimates, based on all of the lines. First, a regression-based Pooled estimate was obtained by using the average of the six means for a given contrast (High, Low, or Divergence) as the data points in a regression. For example, the Pooled High estimate was obtained by regressing the cumulative difference between the mean of all high lines and the mean of all control lines on generations. Second, Falconer also considered a simple Mean estimate for the High, Low, and Divergent comparisons, which was computed as the

604

average of the realized heritability estimates for each comparison over all six replications.

The standard error for the Mean estimate measures the variation seen over the replicates. This is the empirical estimate of the standard error used to compare against the theoretical predictions. The average realized heritability estimates under each of the three different contrasts were very similar: 0.395, 0.331, and 0.369 for High, Low, and Divergent lines (respectively). The three Pooled estimates give very similar estimates. As expected (because OLS estimated standard errors are expected to be downwardly biased), the standard errors for the Mean estimates (based on variation between the replicate lines) were 1.6, 3.3, and 2.3 times larger than the estimated standard error computed using OLS regression of the pooled data.

A second experiment (López-Fanjul and Domínguez 1982; summarized by López-Fanjul 1982) followed 20 replicate Drosophila lines selected for sternopleural bristle number. They found that the empirical standard errors were less than OLS-estimated SEs, which in turn were less than GLS-estimated SEs. This is contrary to the expectation that the OLS-estimated SEs are smaller than the true variances, while the empirical and GLS-estimated SEs are expected to be roughly equal. The authors suggested that this may be at least partly due to a scale effect (LW Chapter 11), as the phenotypic variance greatly decreased during selection.

Finally, (Bohn et al. (1983) examined the variance in response to selection on pupal weight in Trioblium castaneum, examining three replicates from each of six different base populations. The authors found a poor fit when using Hill's (1971) variance formula, which attempts to correct for the effects of selection. Interestingly, a strong correlation was observed between departures in the predicted variance and departures from normality in the base population. The variance among replicate lines drawn from base populations with increasing (absolute) amounts of kurtosis showed larger departures from the predicted value.

\section{Realized Heritability With Rank Data}
Realized heritability estimates are not limited to normally distributed traits. Indeed, the heritability of the liability underlying a threshold trait (Chapter 14) can be computed using a realized heritability estimator, which compared performance of the offspring from different parental classes (Example 14.5; LW Chapter 25). In a similar fashion, Schwartz and Wearden (1959) considered a realized heritability estimator when the data are reduced to ranks, e.g., pecking order. They considered the case where individuals were divided into high- and low-dominance groups, by first ranking the dominance of individuals (in their case, chickens) within a flock, and then assigning those in the upper half to the high group and those in the lower half to the low group. The selection differential was computed as the difference in the mean rank of high and low parents,  $ a_{h} - a_{l} $. Conversely, the selection response was measured as the difference in the mean rank of progeny from high parents versus low parents,  $ b_{h} - b_{l} $, with

\begin{equation}
\widehat{h}_{r}^{2}=\frac{b_{h}-b_{l}}{a_{h}-a_{l}}
\end{equation}

The authors cleverly noted how this estimator can be related to the Mann-Whitney U statistic (for comparing two sample means based on ranks; Conover 1999), and they used this to develop large-sample confidence intervals for the resulting heritability estimate.

\section{Infinitesimal-model Corrections for Disequilibrium}
Our final comment on estimation is that directional selection is expected to decrease the additive variance (and hence the heritability), due to the generation of negative gametic-phase disequilibrium (Chapters 16 and 24). Hence, realized heritability is depressed relative to the (unselected) base population value (Figure 18.4). Under the infinitesimal model, the

605

\begin{figure}[h]
\centering
\includegraphics[width=0.8\linewidth]{example-image}
\caption{Figure}
\end{figure}

majority of this decrease occurs over the first few generations, after which the heritability is essentially at its equilibrium value,  $ \widetilde{h}^{2} $, where (from Equations 16.13a and 16.13c)

\begin{equation}
\begin{align*}\widetilde h^2=\frac{\gamma}{1+\gamma-h^2}\quad{\rm with}\quad\gamma=\frac{2h^2-1+\sqrt{1+4h^2(1-h^2)\kappa}}{2(1+\kappa)}\end{align*}
\tag{18.21}
\label{eq:18_21}
\end{equation}

Here  $ h^2 $ is the (unselected) base population heritability and  $ \kappa $ is a measure of the reduction in phenotypic variance due to selection (the variance following selection being  $ (1 - \kappa)\sigma_z^2 $, Equation 16.10a). In particular,  $ \kappa = \bar{\imath}\left(\bar{\imath} - z_{[1-p]}\right) $ for truncation selection that saves a fraction,  $ p $, of the population for a normal trait (Table 16.1). One can thus approximately treat the realized heritability as an estimate of the equilibrium heritability  $ \hat{h}_r^2 \simeq \bar{h}^2 $, using Equation 18.21 to numerically solve for the base population heritability  $ (h^2) $, an approach first used by Atkins and Thompson (1986). Figure 18.4 shows the relative percentage by which the base population heritability is underestimated by the uncorrected realized heritability.

\begin{table}[h]
\centering
\begin{tabular}{|c|c|}
\hline
Cell 1 & Cell 2 \\
\hline
Cell 3 & Cell 4 \\
\hline
\end{tabular}
\caption{Table Placeholder}
\end{table}

606

If we assume that the infinitesimal model is a reasonable approximation over the course of this experiment, all three methods underestimated the base population heritability by about 13% (see Figure 18.4).

\section{EXPERIMENTAL EVALUATION OF THE BREEDER'S EQUATION}
Although the breeder’s equation requires a number of assumptions (Table 13.2) and, strictly speaking, holds only for a single generation of selection from an unselected base population, the general claim is that it is usually satisfactory (using the base population  $ h^{2} $ value) over a few (5–10) generations (e.g., Falconer 1981). A slightly more refined statement is that the number of generations should not exceed  $ N_{e}/2 $ (Hill 1977a), as drift significantly changes the genetic variance within a line after this amount of time has elapsed (Equation 11.2). After a sufficient number of generations, drift and selection change the genetic variances substantially from their base population values, and the breeder’s equation (using the initial heritability value) fails (Chapters 24–26).

Our concern here is the adequacy of the breeder's equation for short-term artificial selection experiments, with other empirical trends from selection experiments (both short-and long-term) reviewed in Chapters 25 and 26 and with response in natural populations examined in Chapter 20. Starting with the first formal comparisons by Reeve and Robertson (1953) and Clayton and Robertson (1957), a number of authors have compared their observed short-term responses with those predicted from the breeder's equation. As summarized below, the results are mixed. One complication is that most analyses use ordinary (unweighted) least squares, resulting in significantly underestimated standard errors. Likewise, almost all realized heritability estimates have not been adjusted for the expected decline due to gametic-phase disequilibrium.

\section{Most Traits Respond to Selection}
Before delving into the goodness-of-fit of the breeder's equation, it is important to emphasize a fundamental observation in quantitative genetics: Most traits in outbred populations respond to selection. This is such a fundamental observation that the few exceptions are classic and well-known. The first such exception is sex ratio in chickens. Despite enormous economic incentive to create female-biased lines, to our knowledge this has not happened.

The second exception is selection for directional asymmetries, DA, (i.e., “handedness”) in traits in Drosophila (Lewontin 1974). While selection to change the variance in asymmetry is usually successful (Chapter 17), attempts to select for a consistently higher trait value on one specific side of an organism (such as more bristles on the right side) when DA are absent in the base population have historically been unsuccessful (Maynard Smith and Sondhi 1960; Purnell and Thompson 1973; Coyne 1987; Tuinstra et al. 1990; Carter et al. 2009). Because there are numerous examples, such as handedness of the larger claw in fiddler crabs and wing features in many insects (reviewed by Pélabon and Hansen 2008), of such directional asymmetries in nature, this lack of response is intriguing. Two studies offer some potential insight. First, Pélabon et al. (2005) observed changes in DA as a correlated response to selection on an index of wing traits in Drosophila. The key difference from previous unsuccessful experiments is that this index itself showed DA in the control population, while the previously cited failures of response in DA occurred in traits that initially were symmetric. Hence, when DA are already present, it may be possible to amplify them via selection. Second, Rego et al. (2006) found significant genetic variation in DA following hybridization of two closely related Drosophila species, showing that there is at least between-population genetic variation in DA for some traits.

When we move from univariate to multivariate traits, the notion that each trait usually responds to selection is less clear. There often appear to be serious constraints (i.e., little

607

\begin{table}[h]
\centering
\begin{tabular}{|c|c|}
\hline
Cell 1 & Cell 2 \\
\hline
Cell 3 & Cell 4 \\
\hline
\end{tabular}
\caption{Table Placeholder}
\end{table}

genetic variation) for specific combinations of traits, despite significant heritabilities in each component. In such settings, selection on only a specific trait results in a response, while selection on that trait as part of a multivariate selection index may not yield any response, or even may yield a response in the opposite direction of selection on that trait (Volume 3).

\section{Sheridan's Analysis}
One of the most extensive reviews of the fit of the breeder's equation is that of Sheridan (1988), who examined 198 experiments involving laboratory and domesticated animals and compared realized heritabilities with estimates of heritability based on resemblances between relatives. Sheridan first considered those experiments in which the base populations already had an extensive past history of artificial selection on the target trait. In these populations, the response was very poorly predicted by the breeder's equation, with all 11 experiments (6 in Drosophila, 2 in Tribolium and 3 in mice) showing greater than 50% disagreement between the realized and estimated (i.e., base population) heritabilities.

Table 18.1 shows the fit for the remaining 187 experiments, namely, for those traits without an extensive past history of artificial selection. As the data show, the fit was rather poor in many experiments—roughly half have a disagreement of at least 30%, and one in three exceeds 50%. In addition to the biological reasons listed in Table 13.2, there are also design issues that could account for the apparently poor fit. First, none of the experiments reviewed by Sheridan corrected for the expected decline in the realized heritability due to gametic-phase disequilibrium. Second, small absolute disagreements can translate into large relative percentages of disagreement for traits with low heritabilities (Hill and Caballero 1992). For example,  $ |\widehat{h}_r^2 - \widehat{h}^2| = 0.04 $ is a 20% relative disagreement if  $ \widehat{h}_r^2 = 0.2 $, but an 80% disagreement if  $ \widehat{h}_r^2 = 0.05 $. Third, variance in response will also generate some level of disagreement. This raises the questions as to which differences in the table are significant. As a rough approximation, Sheridan considered the disagreement to be significant if it exceeded two standard errors (which amounts to  $ p < 0.05 $ if the data are normally distributed). Assuming the two estimates are uncorrelated, the variance for their difference is the sum of the variance for each estimate, giving a standard error of

\begin{equation}
\mathrm{SE}\left(\widehat{h}^{2}-\widehat{h}_{r}^{2}\right)=\sqrt{\left(\mathrm{SE}[\widehat{h}^{2}]\right)^{2}+\left(\mathrm{SE}[\widehat{h}_{r}^{2}]\right)^{2}}
\tag{18.22}
\label{eq:18_22}
\end{equation}

608

\begin{table}[h]
\centering
\begin{tabular}{|c|c|}
\hline
Cell 1 & Cell 2 \\
\hline
Cell 3 & Cell 4 \\
\hline
\end{tabular}
\caption{Table Placeholder}
\end{table}

\begin{table}[h]
\centering
\begin{tabular}{|c|c|}
\hline
Cell 1 & Cell 2 \\
\hline
Cell 3 & Cell 4 \\
\hline
\end{tabular}
\caption{Table Placeholder}
\end{table}

Taking those 131 experiments with standard errors for both  $ \widehat{h}_{r}^{2} $ and  $ \widehat{h}^{2} $, Table 18.2 shows that 25% had realized heritabilities significantly different from  $ \widehat{h}^{2} $. One problem with this approach is that most of the reported standard errors for  $ \widehat{h}_{r}^{2} $ used by Sheridan were based on unweighted least-squares (OLS), which underestimates the true standard error, giving confidence intervals that are too narrow and thereby inflating the level of significance. The overall 25% of experiments reported by Sheridan as showing a significant departure from the predicted response is thus an upper bound, suggesting that the breeder's equation may not be doing such a poor job after all.

Finally, Sheridan looked at the goodness-of-fit as a function of the duration of the experiment (Table 18.3). Surprisingly, longer experiments tended to have a better fit. While this is contrary to expectations, it could simply be a design artifact. In many cases, longer experiments employ larger population sizes than experiments of shorter duration, thus reducing the effects of drift.

\section{Realized Heritabilities, Selection Intensity, and Inbreeding}
In addition to quantitative differences in the predictions of the breeder's equation discussed above, there are also reported cases of major qualitative departures. For example, the breeder's equation predicts that while selection response should increase with selection intensity, the ratio of response to selection differential, R/S, should be constant. Some studies have reported a dependence of realized heritability on the selection intensity, although a survey of selection experiments found no consistent pattern (Table 18.4).

609

\begin{table}[h]
\centering
\begin{tabular}{|c|c|}
\hline
Cell 1 & Cell 2 \\
\hline
Cell 3 & Cell 4 \\
\hline
\end{tabular}
\caption{Table Placeholder}
\end{table}

\begin{table}[h]
\centering
\begin{tabular}{|c|c|}
\hline
Cell 1 & Cell 2 \\
\hline
Cell 3 & Cell 4 \\
\hline
\end{tabular}
\caption{Table Placeholder}
\end{table}

There are several reasons why a dependency may arise between realized heritability and selection intensity. First, increasing the selection intensity increases gametic-phase disequilibrium, thus reducing  $ \sigma_A^2 $ and hence  $ \widehat{h}_r^2 $ (Chapter 16). As Figure 18.4 shows, the prediction is that (uncorrected) realized heritabilities should decrease with increasing selection intensity (albeit only slightly in many cases). Again, essentially none of the reported selection experiments correct for this expected reduction. Second, allele frequencies are expected to change more rapidly as the strength of selection increases. Whether this results in an increased or decreased response depends on the initial distribution of allele frequencies and effects (see the discussion on genetic asymmetries below).

Finally,  $ N_{e} $ decreases as selection intensity increases (Chapter 3 and 26), thus increasing the amount of inbreeding and (generally) reducing the additive variance (Chapter 11). When inbreeding occurs, the short-term response is generally found to be less than that predicted from the breeder's equation when using variance components estimated from the base population (Table 18.5). This is expected, as inbreeding (generally) decreases the additive genetic variance within a line (Chapter 11), thus reducing response. Exceptions can occur if nonadditive genetic variance is present, in which case inbreeding may actually result in an increase in the within-line additive variance (Chapter 11), increasing response. The effects of drift on long-term experiments are examined in detail in Chapter 26. Inbreeding depression is another complication, which we will address shortly.

610

\begin{figure}[h]
\centering
\includegraphics[width=0.8\linewidth]{example-image}
\caption{Figure}
\end{figure}

\begin{figure}[h]
\centering
\includegraphics[width=0.8\linewidth]{example-image}
\caption{Figure}
\end{figure}

\begin{figure}[h]
\centering
\includegraphics[width=0.8\linewidth]{example-image}
\caption{Figure}
\end{figure}

611

\begin{table}[h]
\centering
\begin{tabular}{|c|c|}
\hline
Cell 1 & Cell 2 \\
\hline
Cell 3 & Cell 4 \\
\hline
\end{tabular}
\caption{Table Placeholder}
\end{table}

\section{Asymmetric Selection Response}
A common design is to perform a divergent selection experiment, wherein replicate lines are selected in opposite directions. Many such experiments (e.g., Figures 18.5–18.8) show different amounts of response in the up versus down direction, a phenomenon referred to as an asymmetric selection response (Falconer 1954). This is in sharp contrast with the expectation from the breeder's equation, which predicts that the absolute magnitude of response should depend only on the absolute value of S, and not on its sign.

There are a variety of possible explanations for asymmetric responses (Table 18.6). They may simply be artifacts of experimental design or analysis. In particular, the prediction of equal positive and negative slopes holds only for plots of cumulative responses versus cumulative selection differentials,  $ R_{C}(t) $ versus  $ S_{C}(t) $. Asymmetry in response based on differences in slope of cumulative response versus generations of selection,  $ R_{C}(t) $ versus t, can thus be misleading, as the different lines may have experienced different amounts of selection. Likewise, differences in response could simply be scale effects (Figure 18.6; see LW Chapter 11). For example, if the genetic variance increases with the mean, the heritability can also increase with the mean, resulting in a faster response in the upwardly selected lines. Asymmetric differences in selection response can also be due to biological features that were not considered by the design. For example, even if the amount of artificial selection is the same in both directions, there may be major differences in natural selection (e.g., up-selected lines might experience lower fertility). Using effective selection differentials (Equation 13.9) corrects for this source of bias, but the investigator often lacks the data (e.g., fertilities for each parent) necessary to compute them. Likewise, we have seen that if a population has been under previous selection, its mean can change even after selection has stopped (Chapter 15). If we start a divergent selection experiment using such a nonequilibrium population as our base, we can bias the response in at least one direction, thus generating an asymmetric response even when the true genetic trend from the experiment is symmetric.

Even though lines may have quite different values of  $ h_{r}^{2} $, there remains the issue of whether these differences are statistically significant. Genetic drift can generate considerable variation between replicate lines, and it is important to distinguish between real differences in the expected response and the variation among realizations of a process with the same (absolute) expected value. Directional trends in the environment can also produce asymmetries, inflating the apparent response in the direction of the trend and retarding it in the opposite direction. An example of this can be seen in the bottom graph in Figure 18.7. If one simply compares the High and Low lines, a clear asymmetric response is seen.

612

\begin{figure}[h]
\centering
\includegraphics[width=0.8\linewidth]{example-image}
\caption{Figure}
\end{figure}

\begin{figure}[h]
\centering
\includegraphics[width=0.8\linewidth]{example-image}
\caption{Figure}
\end{figure}

However, when comparing each line as measured as a deviation from the Control line, the response is much more symmetric.

Although spurious asymmetric responses can result from these defects in design and analysis, true asymmetric responses can be generated by a variety of genetic situations. For example, if the parent-offspring regression is nonlinear, S will not be sufficient to predict the response (Gimelfarb and Willis 1994; Chapters 6, 13, and 24), and it is not surprising that asymmetric responses can be generated in such cases. Failure to initially detect departures from linearity in the base population does not rule out nonlinearity as an explanation for an observed asymmetric response, as the range of variation in the base population may not be sufficient to detect any such departures at the extreme ends of the initial range of phenotypes. As the selected lines diverge, however, differences in the tails of the initial

613

phenotypic distribution can be revealed. Likewise, as previously mentioned, heritability can be a function of the part of the phenotypic distribution from which individuals are drawn (e.g., Skibinski and Shereif 1989).

Characters displaying inbreeding depression (LW Chapter 10) show an asymmetric selection response, with the change in mean from inbreeding accentuating the response in one direction and retarding it in the other. A simple test for inbreeding as an explanation of an asymmetric response is to see whether the mean of an unselected control population (when inbred) changes in the direction of greater response. The effects of inbreeding depression can be corrected by inbreeding a control population to the same level as the divergent selection lines and using the contrast between selected and control lines to estimate response. However, this process is not necessarily as straightforward as it appears, as selection generally increases the amount of inbreeding within a line by decreasing its effective population size (Chapters 3 and 26). Simply keeping the control line at the same size as the selected lines thus underestimates the amount of inbreeding (and hence any correction for inbreeding depression), especially when selection is intense. An alternative tests is to cross two replicate lines, with heterosis in the resulting  $ F_{1} $ indicating inbreeding depression.

The assumption of negligible allele-frequency change over a few generations of selection can be violated when alleles of major effect are segregating. In this setting, asymmetries can arise as a consequence of allele frequencies changing in different directions in the differentially selected populations. An increase in an allele from its initial frequency nearly always results in a different additive variance from that produced by an equal decrease in frequency. Falconer (1954) refers to this feature as genetic asymmetry. This is most easily seen by considering LW Figure 4.6, which shows additive genetic variation as a function of allele frequency for a single diallelic locus. If alleles are completely additive, then  $ \sigma_A^2 $ as a function of allele frequency is symmetric about  $ p = 1/2 $. Suppose that the frequency of an allele that increases the character value is initially below  $ 1/2 $. In upwardly selected lines,  $ \sigma_A^2 $ (and  $ h^2 $) increase as this allele increases to a frequency of 0.5 and decrease when it exceeds this value. Conversely, in downwardly selected lines, the contribution to  $ \sigma_A^2 $ from this locus always decreases. If dominance is present, then  $ \sigma_A^2 $ will no longer be a symmetric function of allelic frequencies (LW Figure 4.8) and asymmetric changes in the contribution to  $ \sigma_A^2 $ from a single locus are almost always expected.

Frankham and Nurthen (1981) provided an interesting example of this phenomena. They started with a base population with the major recessive allele  $ sm^{lab} $ (which greatly reduces abdominal bristle number in Drosophila melanogaster) at low frequency. As shown in Figure 18.8, in two of three down-selected lines, this allele increased in frequency, resulting in a large increase in  $ h^{2} $ as the  $ sm^{lab} $ allele reached intermediate frequencies. Presumably, this allele was either lost or not sampled when Line 1 was formed, as it does not show this pattern. Heritability returned to the base population value as this allele approached fixation, suggesting much smaller allele-frequency change in the residual polygenic genetic variation that remains once the major allele has been removed.

While we have focused on the effects of a single major gene, the effects of unequal allele frequencies apply to any QTL. Alleles of small effect are expected to have slower changes in allele frequencies (and thus are expected to have smaller effects on asymmetry) over the time scales of most short-term experiments. Hence, if many loci of small effect underlie a character, the effects of genetic asymmetry on selection response are expected to be slight unless the number of generations is large. Further, at least some cancellation is expected between the asymmetries in  $ \sigma_{A}^{2} $ generated by allele-frequency changes at different loci. Analogous to directional dominance being required for inbreeding depression (a systematic trend in the direction of dominance over loci; LW Chapter 10), there must also be some systematic trend in the allele frequencies at loci underlying the trait for genetic asymmetries to occur (in the absence of major genes). For example, if most alleles that decrease the trait value tend to be rare, then up-selected versus down-selected lines are expected to show an asymmetric response as allele-frequency changes become significant, with the rare alleles increasing (and increasing variance) in the down-selected, but not the up-selected, lines.

614

\begin{figure}[h]
\centering
\includegraphics[width=0.8\linewidth]{example-image}
\caption{Figure}
\end{figure}

\begin{figure}[h]
\centering
\includegraphics[width=0.8\linewidth]{example-image}
\caption{Figure}
\end{figure}

When might such an asymmetric distribution of allele frequencies be present? One situation is that where there has been a recent history of selection on the trait before the start of the

615

\begin{figure}[h]
\centering
\includegraphics[width=0.8\linewidth]{example-image}
\caption{Figure}
\end{figure}

After selection:
only clone A survives

\begin{figure}[h]
\centering
\includegraphics[width=0.8\linewidth]{example-image}
\caption{Figure}
\end{figure}

artificial selection experiment.

Components of reproductive fitness, such as fecundity and development time, are expected to have asymmetric allele frequencies, as natural selection increases the frequency of alleles increasing fitness. This condition suggests that asymmetric responses in the selection for reproductive traits are expected, and it further predicts that response should be larger in lines selected for a decrease in reproductive fitness. This was seen in a review by Frankham (1990), who found a larger response in the direction of reduced reproductive fitness in 24 of 30 experiments. Frankham suggested that the presence of rare recessives decreasing reproductive fitness was the most likely cause for this trend. As Figure 25.2 will show, very marked asymmetric responses are expected in such situations.

\section{Reversed Response}
The most extreme departure from the breeder's equation is a reversed response, a response in the opposite direction of selection. Negative maternal effects can result in such a response (Figures 15.3 and 15.4). Likewise, they can arise from sufficiently large genotypic-dependent environmental variances (Haldane 1931; Chapter 17). Haldane's intuition was based on selection with two asexual clones, the first with a higher mean but smaller environmental variance than the second. As shown in Figure 18.9, if the most extreme individuals are selected, there will be an excess of the clone with the smaller mean but higher variance, resulting in a decrease in mean value in the next generation. Wright (1969) expanded Haldane's model to a single diallelic locus. Gimelfarb (1986b) presented a particularly interesting analysis when there is a multiplicative genotype × environment interaction, showing for this model that

616

while phenotypes are subjected to directional selection, the nature of the interaction between genotype and environment is such that genotypic values can actually be under either stabilizing or disruptive selection. The interesting feature of both Haldane's and Gimelfarb's models is that reversed response is most probable when selection is very intense, as this is the setting most favorable to genotypes with high variances (Chapter 17).

Important, and not yet widely appreciated, sources of reversed response are associative effects (also called social effects or indirect genetic effects), of which maternal effects are a special case. As introduced by our discussions on maternal effects in Chapter 15, and more fully developed in Chapter 22, the environment that a focal individual experiences can be strongly influenced by its neighbors. In such cases, an individual's phenotype is the result of both direct effects, due to its own intrinsic genes, and social effects, from its neighbors. If the latter have a heritable component, these effects can also evolve. If direct and social effects are genetically negatively correlated, a positive response in the direct trait can be overcome by a negative social response, resulting in a net decrease in trait value. Consider selection for egg production in caged chickens. Aggressive individuals gather more food in the cage thus producing more eggs. With selection based solely on individual egg production, more aggressive chickens are chosen (i.e., those with negative associative effects), creating a deleterious cage environment that can overshadow any direct genetic advance in terms of number of eggs, resulting in a reduction in the number of eggs per cage in the next generation.

\section{CONTROL POPULATIONS AND EXPERIMENTAL DESIGNS}
We conclude this chapter by considering several specialized issues in the analysis of selection experiments, topics that may be skipped by the casual reader. We first examine the strengths and weaknesses of different types of designs (such as the use of a control population), then consider optimal designs, and finish with a discussion of a modified line-cross analysis (LW Chapter 9) to examine the genetic components of a selection response.

One complication with estimating realized heritabilities is distinguishing between genetic versus environmental trends. For example, Newman et al. (1973) found that 60% of the increase in yearling weight in a selected line of Shorthorn cattle was due to environmental, rather than genetic, improvement. Using a large unselected control population reared under the same environmental conditions as the selected line or lines allows for correction of such between-generation environmental changes (e.g., the bottom graph in Figure 18.7). Hill (1972a, 1972b) reviewed this approach and some of the features seen in control populations from a number of different species. The use of control populations is not a foolproof approach for removing environmental trends, as the control and selected lines can evolve different genotype × environment interactions. Likewise, extensive genetic drift in the mean of the control population can also result in biases, resulting in the corrected (selection minus control) values underestimating or overestimating the true selection trend. When full pedigree data are available, the powerful methods of mixed-model analysis (Chapter 19) can be used to remove genetic and environmental trends, even in the absence of control populations (although incorporating such controls is preferred). Undetected environmental trends are especially problematic in the analysis of response in natural populations (Chapter 20).

\section{Basic Theory of Control Populations}
Assuming no genotype × environment interaction, the true mean of a line in generation t can be decomposed as  $ \mu + g_t + b_t $, where  $ \mu $ is the base population mean,  $ g_t $ is the change in mean breeding value due to selection and drift, and  $ b_t $ is the change in mean due to environmental change. Under this model, observed means,  $ \overline{z} $, for a selected (s) and control (c) population reared in a common environment can be decomposed as

\begin{equation}
\overline{z}_{s,t}=\mu+g_{s,t}+b_{s,t}+e_{s,t}
\tag{18.23a}
\label{eq:18_23a}
\end{equation}

\begin{equation}
\overline{z}_{c,t}=\mu+g_{c,t}+b_{c,t}+e_{c,t}
\tag{18.23b}
\label{eq:18_23b}
\end{equation}

617

where  $ e_t $ is the error in estimating the mean breeding value  $ (\mu + g_t) $ from the observed mean, corrected for the change in the environment  $ (\overline{z}_t - b_t) $, and has an expected value of 0.

Assuming that the breeder's equation holds, the expected total (i.e., cumulative) response at generation t is  $ E[R_C(t)] = E[g_{s,t}] = h^2 S_C(t) $. Under drift alone, there is no expected directional change in the mean breeding value of the control population,  $ E[g_{c,t}] = 0 $. If we assume that there are no genotype × environment interactions (i.e.,  $ b_{s,t} = b_{c,t} $), the contrast between the selected and control populations has an expected value of

\begin{equation}
E[\overline{z}_{s,t}-\overline{z}_{c,t}]=E[g_{s,t}]-E[g_{c,t}]=h^{2}S_{C}(t)
\tag{18.24a}
\label{eq:18_24a}
\end{equation}

If the control population is small,  $ g_{c,t} $ can drift significantly away from zero, resulting in an overestimation (or underestimation) of the true selection response. If the goal is simply to remove an environmental trend, control populations should be kept as large as possible so that  $ g_{c,t} $ is close to zero. However, as previously mentioned, control populations are also used in an attempt to correct for any effects of inbreeding depression. In such cases, the size of the control population is usually on par with that for the selected population, meaning that significant drift has likely occurred. In this setting, the use of a control population may introduce additional uncertainty in the estimate of the response (which may be more than compensated for by the control accounting for inbreeding depression and environmental trends).

If genotype  $ \times $ environmental interactions are present, then

\begin{equation}
E[\overline{z}_{s,t}-\overline{z}_{c,t}]=h^{2}S_{C}(t)+(b_{s,t}-b_{c,t})
\tag{18.24b}
\label{eq:18_24b}
\end{equation}

resulting in  $ \overline{z}_{s,t} - \overline{z}_{c,t} $ being a potentially biased estimator of  $ h^{2}S_{C}(t) $. In principle, if the environmental trends are positively correlated between populations, then the use of a control will still improve the estimate of the genetic trend. If they are negatively correlated, however, the use of a control can lead to a less accurate estimate than simply using a selected population without a control. As discussed in Chapter 19, mixed-model analyses may be able to provide some insights (as they estimate the  $ b_{x,t} $), but they require extensive information (full pedigrees) and that the model assumptions hold.

If a control population is used, the (per-generation,  $ R_{t} $, and cumulative,  $ R_{C}[t] $) responses and differentials are estimated by

\begin{equation}
R_{t}=\left(\overline{z}_{s,t}-\overline{z}_{c,t}\right)-\left(\overline{z}_{s,t-1}-\overline{z}_{c,t-1}\right)
\tag{18.23a}
\label{eq:18_23a_2}
\end{equation}

\begin{equation}
R_{C}(t)=\overline{z}_{s,t}-\overline{z}_{c,t}
\tag{18.25b}
\label{eq:18_25b}
\end{equation}

\begin{equation}
S_{t}=\overline{z}_{s,t-1}^{*}-\overline{z}_{s,t-1}
\tag{18.25c}
\label{eq:18_25c}
\end{equation}

\begin{equation}
S_{C}(t)=\sum_{i=1}^{t}S_{i}
\tag{18.25d}
\label{eq:18_25d}
\end{equation}

where  $ \overline{z}^{*} $ denotes the mean of the selected individuals. If selection differs between the sexes

\begin{equation}
S_{t}=\frac{S_{t}(m)+S_{t}(f)}{2}
\end{equation}

where  $ S_{t}(m) $ and  $ S_{t}(f) $ are, respectively, the observed selection differentials on males and females.

Muir (1986a, 1986b) has suggested that an analysis of covariance approach (along the lines of LW Equation 10.12) be used when extensive G × E is expected. In its simplest form, Muir's idea is to adjust the mean of the selected population in generation i by using the deviation of  $ \overline{z}_{c,i} $ from its grand mean over the course of the experiment,  $ \overline{z}_{c,\cdot} $, namely

\begin{equation}
\overline{{z}}_{s,i}^{\prime}=\overline{{z}}_{s,i}-\beta\big(\overline{{z}}_{c,i}-\overline{{z}}_{c,.}\big)
\end{equation}

If there is no genotype × environment interaction (so that the environmental effect at time i has the same influence on both the selected and control lines), then  $ \beta = 1 $ and we recover

618

our previous results. However, if there is  $ G \times E $, such that the environmental deviation influences the control differently from the selected population, then this expression provides a less biased correction than simply subtracting off the mean of the control. Note that this covariate approach at least partly corrects for the most extreme case, namely that in which the environmental effect results in different signs in the control and selective populations ( $ \beta < 0 $).

\section{Divergent Selection Designs}
A related approach is the divergent (or bidirectional) selection design, wherein one compares lines that were selected in opposite directions. Again assuming there are no significant genotype × environment interactions between lines, the basic statistical model for this design is

\begin{equation}
\overline{z}_{u,t}=\mu+g_{u,t}+b_{u,t}+e_{u,t}
\tag{18.26a}
\label{eq:18_26a}
\end{equation}

\begin{equation}
\overline{z}_{d,t}=\mu+g_{d,t}+b_{d,t}+e_{d,t}
\tag{18.26b}
\label{eq:18_26b}
\end{equation}

where the subscripts u and d refer to the upwardly and downwardly selected lines, respectively. Using this design (and assuming that  $ b_{u,t}=b_{d,t} $), the responses and differentials are estimated by

\begin{equation}
R_{t}=\left(\overline{z}_{u,t}-\overline{z}_{u,t-1}\right)-\left(\overline{z}_{d,t}-\overline{z}_{d,t-1}\right)
\tag{18.27a}
\label{eq:18_27a}
\end{equation}

\begin{equation}
R_{C}(t)=\overline{z}_{u,t}-\overline{z}_{d,t}
\tag{18.27b}
\label{eq:18_27b}
\end{equation}

\begin{equation}
S_{t}=\left(\overline{z}_{u,t-1}^{*}-\overline{z}_{u,t-1}\right)-\left(\overline{z}_{d,t-1}^{*}-\overline{z}_{d,t-1}\right)
\tag{18.27c}
\label{eq:18_27c}
\end{equation}

Again, the expected response (using Equations 18.27a and 18.27c for the response and selection differentials) is simply  $ R = h^{2}S $. More generally, if  $ S_{x,i} $ denotes the selection differential in line x at generation i, the total (cumulative) expected divergence between lines is

\begin{equation}
R_{C}(t)=h^{2}\sum_{i=1}^{t}\left(S_{u,i}-S_{d,i}\right)=h^{2}\left(S_{C,u}-S_{C,d}\right)
\tag{18.27d}
\label{eq:18_27d}
\end{equation}

\section{Variance in Response}
Recall the pure-drift approximation for the variance (Equation 18.8) and the within-line covariance across generations (Equation 18.9) for the design of a single selected line. Here we present corresponding expressions for the selection plus control and divergence selection designs (assuming there are no genotype × environmental interactions between lines). For unidirectional selection plus a control population

\begin{equation}
\begin{aligned}R_{C}(t)=\overline{z}_{s,t}-\overline{z}_{c,t}&=\left(\mu+g_{s,t}+b_{t}+e_{s,t}\right)-\left(\mu+g_{c,t}+b_{t}+e_{c,t}\right)\\&=\left(g_{s,t}-g_{c,t}\right)+\left(e_{s,t}-e_{c,t}\right)\end{aligned}
\tag{18.28}
\label{eq:18_28}
\end{equation}

Similarly, for divergent selection

\begin{equation}
R_{C}(t)=\overline{z}_{s u,t}-\overline{z}_{s d,t}=\left(g_{s u,t}-g_{s d,t}\right)+\left(e_{s u,t}-e_{s d,t}\right)
\tag{18.29}
\label{eq:18_29}
\end{equation}

Because each term in Equations 18.28 and 18.29 is independent, applying Equations 18.4 and 18.5 yields a pure-drift approximation for the variance in response in generation t of

\begin{equation}
\sigma^{2}\left[R_{C}(t)\right]=\left(2f_{t}+B_{0}\right)h^{2}\sigma_{z}^{2}+B_{t}\sigma_{z}^{2}\simeq\left(t A+B_{0}\right)h^{2}\sigma_{z}^{2}+B_{t}\sigma_{z}^{2}
\tag{18.30a}
\label{eq:18_30a}
\end{equation}

and the covariance between generations within the same line

\begin{equation}
\begin{aligned}\sigma\left[R_{C}(t),R_{C}(t^{\prime})\right]&=\left(2f_{t}+B_{0}\right)h^{2}\sigma_{z}^{2}\\&\simeq\left(t A+B_{0}\right)\sigma_{z}^{2}h^{2}\qquad\text{for}t<t^{\prime}\end{aligned}
\tag{18.30b}
\label{eq:18_30b}
\end{equation}

619

Table 18.7 Coefficients for the pure-drift variances and covariances in response (Equations 18.30a and 18.30b).  $ M_{x,t} $ individuals are sampled in population x at generation t, of which  $ N_x $ are allowed to reproduce. The subscripts x = s and c refer to the selected and control populations, x = u and d refer to the up-selected and down-selected lines, and  $ f_{x,t} $ refers to the amount of inbreeding in population x at time t.

Selection in a single direction without a control line. Equation 18.30a has an extra term,  $ \sigma_{b}^{2} $, accounting for the between-generation variation in environmental effects.

\begin{equation}
f_{t}=f_{s,t},\qquad A=\frac{1}{N_{s}},\qquad B_{t}=\frac{1}{M_{s,t}}\quad for t\geq0
\end{equation}

Selection in a single direction with a control line.

\begin{equation}
f_{t}=f_{s,t}+f_{c,t},\qquad A=\frac{1}{N_{s}}+\frac{1}{N_{c}},\qquad B_{t}=\frac{1}{M_{s,t}}+\frac{1}{M_{c,t}}\quad for t\geq0
\end{equation}

Divergent selection without a control line.

\begin{equation}
f_{t}=f_{u,t}+f_{d,t},\qquad A=\frac{1}{N_{u}}+\frac{1}{N_{d}},\qquad B_{t}=\frac{1}{M_{u,t}}+\frac{1}{M_{d,t}}\quad for t\geq0
\end{equation}

where the (design-specific) coefficients, A and  $ B_t $, are given in Table 18.7. If the number of reproducing individuals varies over time, then the tA term is replaced by the sum over time of the components in A, for example by  $ \sum_{i=1}^t (1/N_i) $, or related expressions. Finally, recall (Equation 18.3) that for unidirectional selection without a control, the variance in response has an additional term,  $ \sigma_b^2 $, accounting for the between-generation environmental variation.

\section{Control Populations and Variance in Response}
When does using a control population in a undirectional selection experiment reduce the variance in response? Equation 18.30a, along with the coefficients from Table 18.7, gives the expected variance with and without the use of a control. Assuming  $ M = M_s = M_c $ and  $ N = N_s = N_c $, this equation yields

\begin{equation}
\sigma^{2}\left[R_{C}(t)\right]_{control}-\sigma^{2}\left[R_{C}(t)\right]_{no control}=\left(\frac{t}{N}+\frac{1}{M_{0}}\right)h^{2}\sigma_{z}^{2}+\frac{1}{M}\sigma_{z}^{2}-\sigma_{b}^{2}
\tag{18.31a}
\label{eq:18_31a}
\end{equation}

Assuming the between-line drift variance dominates (terms involving M can be ignored), the condition for the variance in response with a control to be larger than the response without one is approximately

\begin{equation}
t\sigma_{z}^{2}h^{2}/N>\sigma_{b}^{2}
\tag{18.31b}
\label{eq:18_31b}
\end{equation}

Hence, regardless of the value of  $ \sigma_{b}^{2} $, if sufficient generations are used, the optimal design (in terms of giving the smallest expected variance in response) is not to use a control. However, this approach runs the risk that an undetected directional environmental trend will

620

Table 18.8 Coefficients of variation (CV) for various designs, assuming the pure-drift approximation and that  $ \sigma^2[R_C(t)] \simeq t A h^2 \sigma_z^2 $. The latter assumes that the selection experiment is of sufficient duration that the among-line drift dominates (i.e.,  $ t A \gg B_0 $ and  $ t A h^2 \gg B_t $, with these parameters defined in Table 18.7). We assume that the absolute selection intensity in all selected lines is  $ \bar{\nu} $.

Selection in a single direction with a control line.

\begin{equation}
E\left[R_{C}(t)\right]=t h^{2}\bar{\imath}\sigma_{z},\qquad\mathrm{C V}\left[R_{C}(t)\right]\simeq\frac{1}{h\bar{\imath}}\sqrt{\frac{2}{N t}}
\tag{18.32a}
\label{eq:18_32a}
\end{equation}

Selection in a single direction without a control line.

\begin{equation}
E\left[R_{C}(t)\right]=t h^{2}\bar{\imath}\sigma_{z},\qquad\mathrm{C V}\left[R_{C}(t)\right]\simeq\frac{1}{h\bar{\imath}}\sqrt{\frac{1}{N t}}+\frac{1}{t h^{2}\bar{\imath}}\left(\frac{\sigma_{d}}{\sigma_{z}}\right)
\tag{18.32b}
\label{eq:18_32b}
\end{equation}

Divergent selection without a control line.

\begin{equation}
E\left[R_{C}(t)\right]=2t h^{2}\bar{\imath}\sigma_{z},\qquad\mathrm{C V}\left[R_{C}(t)\right]\simeq\frac{1}{h\bar{\imath}}\sqrt{\frac{1}{2N t}}
\tag{18.32c}
\label{eq:18_32c}
\end{equation}

compromise the estimated heritability. Thus, the cost from to an increase in SE when using a control seems very worth enduring in many cases, given the more general benefits of having one or more control lines.

\section{OPTIMAL EXPERIMENTAL DESIGNS}
As Equation 18.31b illustrates, it is not entirely obvious which design is optimal. What in general can we say? The coefficient of variation (CV) of the selection response

\begin{equation}
\mathrm{CV}\left[R_{C}(t)\right]=\frac{\sigma\left[R_{C}(t)\right]}{E\left[R_{C}(t)\right]}
\tag{18.33}
\label{eq:18_33}
\end{equation}

is especially useful in comparing efficiencies of different designs, as it is independent of  $ \sigma_{z}^{2} $. Further, it provides an appropriate measure of comparative efficiencies when the expected response differs between designs. Table 18.8 gives expressions for the CV under some simplifying assumptions. The coefficient of variation is a function of  $ tN $, the total number of adults selected during the course of the experiment, provided drift variance dominates the error variance. A short experiment with many selected adults per generation thus gives the same expected CV as a longer experiment with fewer adults per generation (provided the total numbers,  $ tN $, are the same). However, if the error variance is nontrivial relative to the drift variance (as would be expected if  $ h^{2} $ is small), increasing the duration of the experiment results in some improvement in precision (Hill 1980).

As an example of using the CV of response, consider unidirectional selection without a control population versus divergent selection. Which is more efficient if the same total number of adults are selected (e.g.,  $ N $ under unidirectional selection and  $ N_d = N_u = N/2 $ under divergent selection)? If there is no between-generation environmental variance, then both designs are equally efficient, while divergent selection is more efficient if  $ \sigma_b^2 > 0 $.

Example 18.7. Suppose we plan to select the upper 5% of a population for a normally distributed character with  $ h^2 = 0.25 $. What value of  $ Nt $ is needed for the expected CV of response to be no greater than 0.01 if no control population is used? From Example 14.1,  $ E[\bar{t}] = 2.06 $ if the population is large, and slightly less in small populations (for simplicity we

621

assume the large population value applies). To solve for the value of  $ N_t $, we use the expression for the CV in response without a control line (Equation 18.32b), and further assume that the drift variance dominates, so that we ignore the contribution from  $ \sigma_b^2 $. This yields

\begin{equation}
CV=0.01=\frac{1}{0.5\times2.06}\times\sqrt{\frac{1}{Nt}}
\end{equation}

which has a solution of  $ Nt \simeq 9426 $. Hence, during the entire course of the experiment using a total of at least 9426 selected parents returns an approximate expected CV of less than 1%. If the desired CV is 0.05 or 0.10, then  $ Nt \simeq 377 $ and  $ Nt \simeq 94 $, respectively. It is important to stress that these numbers are somewhat misleading, in that N is the number of saved individuals, while M = N/p is the number of measured individuals. For p = 0.05, the value of Mt is 1/0.05 = 20 times the value of  $ Nt $, so one must measure roughly 18,900, 7540, and 1880 individuals to obtain these desired CVs.

\section{Nicholas' Criterion}
An alternative criterion for choosing  $ Nt $ was suggested by Nicholas (1980). Often the investigator is interested in ensuring that at least a certain response will occur with a preset probability. To a reasonable approximation, the expected mean value in any given replicate selection line after  $ t $ generations of selection is normally distributed, with a mean of  $ E[R_C(t)] $ and a variance of  $ \sigma^2[R_C(t)] $. Consider the probability that the observed response is at least  $ \beta $ of the expected response

\begin{equation}
\begin{aligned}\Pr\left(R_{C}(t)>\beta E[R_{C}(t)]\right)&=\Pr\left(\frac{R_{C}(t)-E[R_{C}(t)]}{\sigma\left[R_{C}(t)\right]}>\frac{(\beta-1)E[R_{C}(t)]}{\sigma\left[R_{C}(t)\right]}\right)\\&=\Pr\left(U>\frac{\beta-1}{\mathbf{CV}\left[R_{C}(t)\right]}\right)\end{aligned}
\tag{18.34}
\label{eq:18_34}
\end{equation}

where $U$ is a unit normal random variable. Note that the probability that the observed response exceeds the expected response ($\beta = 1$) is one half, as $\Pr[U > 0] = 1/2$.

Example 18.8. Again suppose that  $ \bar{\imath} = 2.06 $,  $ h^2 = 0.25 $, and the design is unidirectional selection without a control population. What value of  $ Nt $ is required in order to obtain a 95% probability that the observed response is at least 90% of its expected response? Here,  $ \beta = 0.9 $ and  $ \Pr[U > -1.65] = 0.95 $. Hence,

\begin{equation}
\frac{\beta-1}{CV\left[R_{C}(t)\right]}=\frac{-0.1}{CV\left[R_{C}(t)\right]}=-1.65
\end{equation}

Rearranging yields

\begin{equation}
CV\left[R_{C}(t)\right]=\frac{1}{0.5\times2.06}\sqrt{\frac{1}{Nt}}=\frac{0.1}{1.65}
\end{equation}

implying that  $ Nt \simeq 257 $. As in Example 18.7, the required total number of measured individuals is 1/p times this value, yielding  $ Mt = 257/0.05 \simeq 5140 $.

\section{Replicate Lines}
There is no loss of efficiency when replicate lines are used (Hill 1980). To see this, let  $ \overline{z}_{i,t} $ for  $ 1 \leq i \leq r $ be the sample mean for replicate population  $ i $ at time  $ t $. The overall mean is

\begin{equation}
\overline{z}_{t}=\frac{1}{r}\sum_{i=1}^{r}\overline{z}_{i,t}
\tag{18.35a}
\label{eq:18_35a}
\end{equation}

622

If we take the variances and assume each line is independent

\begin{equation}
\sigma_{\overline{z}}^{2}(t)=\frac{1}{r^{2}}\sum_{i=1}^{r}\sigma_{\overline{z}_{i}}^{2}(t)=\frac{1}{r}\sigma_{\overline{z}_{1}}^{2}(t)
\tag{18.35b}
\label{eq:18_35b}
\end{equation}

If the number sampled and number used as parents within a replicate are  $ M^* = M/r $ and  $ N^* = N/r $, respectively, then it is easily seen from Table 18.8 that the variance of a replicate line is simply r times the variance of a population with N and M. Hence, variance in the sample mean from r replicate lines with  $ N^* $ and  $ M^* $ is the same as the variance with a single line with N and M, provided that the number of individuals within each replicate line is sufficiently large to avoid significant inbreeding. Richardson et al. (1968), Irgang et al. (1985), and Muir (1986a, 1986b) developed regression approaches that correct for between-generation environmental changes when replicate lines are used.

\section{LINE-CROSS ANALYSIS OF SELECTION EXPERIMENTS}
Plant-breeding schemes often involve two (or more) generations for each cycle of selection (Chapter 23; Volume 3). The amount of selective progress can be examined by growing seeds from different cycles together in a common garden thus offering a direct assessment of the genetic response. For example, one could regress the mean of a line on its selection-cycle number (e.g., Burton et al. 1971). An even more powerful design involves growing crosses between remnant seed from different cycles of selection in a common garden. Such a line-cross analysis yields summary statistics about the genetic nature of the divergence between lines, such as the role of additive versus nonadditive effects (LW Chapters 9 and 20). As we saw in LW Chapter 20, a large number of parameters are required when one has a collection of unrelated lines (i.e., the general and specific combining abilities of each line and each pair, respectively). However, in the analysis of a selection experiment, lines are not unrelated but rather are genetically connected, requiring far fewer parameters to model. The key to connecting lines from different cycles is to assume a constant rate of allele-frequency change over the course of the experiment (Hammond and Gardner 1974; Smith 1979a, 1979b, 1983; Melchinger and Flachenecker 2006), which is referred to as a generation means analysis (GMA). With a large number of loci, each of which has a small effect, this assumption may be reasonable. If there are major alleles present or selection runs long enough such that substantial allele-frequency change occurs, then this model becomes more problematic. While commonly used in plant breeding, GMA is applicable to any organism for which an equivalent of “remnant seed” is available, such as frozen breeding stock or lines extracted from different generations of selection.

\section{The Simple Additive Model}
To see the logic behind a GMA, consider the simplest case: two alleles at each locus, with only additive effects and no G × E present. Focusing on a particular locus, the genotypes contribute values of 0 : a : 2a to the overall mean, where we sum contributions across all loci (because we assumed there is no epistasis). If p denotes the frequency of the favorable allele, then the contribution from this locus is 2ap. If we further assume a roughly constant rate of change in allele frequency (at least over a few generations), then the frequency, p(k), of the favorable allele in cycle k of selection is

\begin{equation}
p(k)\simeq p+k\Delta p
\end{equation}

where p is the initial frequency and  $ \Delta p $ is the per-generation rate of change for this locus. Summing over all loci, the mean following k cycles of selection is

\begin{equation}
C_{k}=\mu+\sum_{i}2a_{i}p_{i}(k)\simeq\mu+\sum_{i}2a_{i}(p_{i}+k\Delta p_{i})=\mu+\sum_{i}2a_{i}p_{i}+k\sum_{i}2a_{i}\Delta p_{i}
\end{equation}

623

Next, we define

\begin{equation}
m=\mu+2\sum p_{i}a_{i}\qquadand\qquad A=2\sum\Delta p_{i}a_{i}
\tag{18.36}
\label{eq:18_36}
\end{equation}

where m is the line mean before selection and A is a weighted measure of change in the mean given constant allele-frequency change,  $ \Delta p_{i} $, at a locus (which can vary over loci). The expected line mean from selection cycle k becomes

\begin{equation}
C_{k}=m+A\cdot k
\tag{18.37a}
\label{eq:18_37a}
\end{equation}

Now suppose individuals from cycles k and j are crossed. With the frequencies of the favorable allele in these lines being  $ p(k) $ and  $ p(j) $,  $ [1-p(k)] $  $ [1-p(j)] $ is the chance of getting the 0-valued homozygote, while  $ p(k) $  $ [1-p(j)]+p(j) $  $ [1-p(k)] $ is the chance of getting an a-value heterozygote, and  $ p(k)p(j) $ is the chance of getting a 2a homozygote. The resulting contribution from a particular locus to this cross becomes

\begin{equation}
\begin{align*}0\cdot\left\{\left[1-p(k)\right]\left[1-p(j)\right]\right\}+a\cdot\left\{p(k)\left[1-p(j)\right]+p(j)\left[1-p(k)\right]\right\}+2a\cdot\left\{p(k)p(j)\right\}\\=a\cdot\left[p(k)+p(j)\right]=a\cdot\left[2p+(k+j)\Delta p\right]\end{align*}
\end{equation}

Summing over all loci, and recalling Equation 18.36, we get an expected mean for this cross of

\begin{equation}
C_{k\times j}=m+A\cdot\frac{k+j}{2}
\tag{18.37b}
\label{eq:18_37b}
\end{equation}

Under this simple additive model, the mean of any particular population (a line,  $ C_k $, or the progeny from a line cross,  $ C_{k \times j} $) is a function of the constants m and A and the selection cycle number. The goodness-of-fit can be assessed by comparing the predicted and actual mean values (LW Chapter 9).

Equations 18.37a and 18.37b have a simple form because of the assumption of strict additivity (no dominance or epistasis). However, when loci show dominance, the expressions for line and line-cross means are more complex. Line means may then change under selfing (assuming that there is inbreeding depression), although any such change allows us to estimate certain dominance components. Further, when dominance is present, care is required when considering the progeny of line crosses. We can regard the progeny from a particular cross (say  $ C_{k \times j} $) as an  $ F_1 $, and can generate an  $ F_2 $ by either selfing individuals (when possible) or by letting the  $ F_1 $ randomly mate. Under strict additivity, the means of both types of  $ F_2 $ offspring (from either selfing or random mating) are still given by Equation 18.37b, the mean value of the  $ F_1 $. When loci display dominance, the  $ F_1 $ mean and the two different types of  $ F_2 $ all differ from each other. Again, these differences offer opportunities to estimate various dominance components, as we now will describe.

\section{The Hammond-Gardner Model}
Motivated by the general diallel analysis of Gardner and Eberhart (1966), Hammond and Gardner (1974) extended the simple additive model to allow for dominance, which requires three additional parameters  $ (D_{0}, D, \text{ and } D_{q}) $. Assuming all lines start from the same base population, as derived below in Example 18.10, the Hammond-Gardner (HG) model returns means for various lines and line crosses of:

\begin{equation}
C_{k}=m+A\cdot k+D_{0}+D\cdot k+D_{q}\cdot k^{2}
\tag{18.38a}
\label{eq:18_38a}
\end{equation}

\begin{equation}
C_{k}(s)=m+A\cdot k+\frac{1}{2}\left(D_{0}+D\cdot k+D_{q}\cdot k^{2}\right)
\tag{18.38b}
\label{eq:18_38b}
\end{equation}

\begin{equation}
C_{k\times j}=m+\frac{A}{2}(k+j)+D_{0}+\frac{1}{2}\left[D\cdot(k+j)+2D_{q}\cdot k j\right]
\tag{18.38c}
\label{eq:18_38c}
\end{equation}

\begin{equation}
C_{k\times j}(s)=m+\frac{A}{2}(k+j)+\frac{D_{0}}{2}+\frac{1}{4}\left[D\cdot(k+j)+2D_{q}\cdot k j\right]
\tag{18.38d}
\label{eq:18_38d}
\end{equation}

\begin{equation}
C_{k\times j}(rm)=C_{k\times j}+\frac{D_{q}}{4}(k-j)^{2}
\tag{18.38e}
\label{eq:18_38e}
\end{equation}

624

where  $ (s) $ and  $ (rm) $, respectively, denote the selfed and randomly mated next generations given the cross or line. For example,  $ C_4(s) $ is the mean among the selfed progeny of individuals from cycle four. Likewise, for  $ C_{2\times6}(rm) $, we cross cycle two and cycle six individuals, let their offspring (the  $ F_1 $,  $ C_{2\times6}[rm] $) randomly mate, and then measure their resulting progeny (the  $ F_2 $). As expected, when all of the dominance terms are zero, Equations 18.38a–18.38e reduce to Equations 18.37a and 18.37b. With no dominance, there is no change in the line mean following inbreeding, so  $ C_k = C_k(s) $, and the  $ F_1 $ progeny mean from a cross has the same mean as the progeny from either a selfed or randomly mated  $ F_2 $,  $ C_{k\times i} = C_{k\times i}(s) = C_{k\times i}(rm) $.

Five parameters  $ (m, A, D, D_0, $ and  $ D_q) $ are required under the full model, and closer inspection of the coefficients on these model parameters in Equations 18.38a–18.38e reveals what types of crosses are required to estimate parameters of interest. If the investigator only has the unselfed line means,  $ C_k $, and their  $ F_1 $ crosses,  $ C_{k \times j} $, A and D are not separately estimable, as they appear as  $ k(A + D) $ in expressions for these means (see Equations 18.38a and 18.38c). Obtaining separate estimates requires having different coefficients on the A and D terms, as occurs with either selfed lines or selfed  $ F_2 $ crosses (Equations 18.38b and 18.38d). These particular crosses provide additional information on the genetic nature of the between-generation divergence that is not given by the simple line means, or their  $ F_1 $s, alone. Likewise, one cannot obtain separate estimates of m and  $ D_0 $ unless there is at least one selfed line in the study. These terms enter as  $ m + D_0 $ in all randomly mating populations, but they enter as  $ m + D_0 / 2 $ in selfed lines.

As with the strictly additive model, we can express the parameters of the Hammond-Gardner model in terms of locus-specific parameters, with m and A as defined by Equation 18.36. Letting  $ d_i $ be the dominance for locus i, which now makes the genotypic values 0 :  $ a_i + d_i : 2a_i $, the dominance-related parameters,  $ D_0 $, D, and  $ D_q $, in Equations 18.38a–18.38e are as follows.  $ D_0 $ is the contribution (to the mean) from dominance in the initial (cycle-zero) population

\begin{equation}
D_{0}=2\sum p_{i}(1-p_{i})d_{i}
\tag{18.39a}
\label{eq:18_39a}
\end{equation}

The role of dominance in response in the line means enters through D and  $ D_{q} $, which are the linear and quadratic regression coefficients on cycle number. D is calculated by

\begin{equation}
D=2\sum\Delta p_{i}(1-2p_{i})d_{i}
\tag{18.39b}
\label{eq:18_39b}
\end{equation}

and it appears in Equations 18.38a–18.38e as terms of the form  $ D \cdot k $, namely, D times the appropriate selection cycle number, k. Recalling (LW Equation 4.10b) that the average effect,  $ \alpha $, of an allelic substitution (replacing an unfavorable allele by a favorable one) is given by  $ \alpha = a + d(1 - 2p) $, we have

\begin{equation}
\begin{align*}A+D&=2\left(\sum\Delta p_i a_i+\sum\Delta p_i(1-2p_i)d_i\right)\\&=2\left(\sum\Delta p_i\left[a_i+(1-2p_i)d_i\right]\right)=2\sum\Delta p_i\alpha_i\end{align*}
\tag{18.39c}
\label{eq:18_39c}
\end{equation}

with  $ A+D $ representing the per-generation change in the trait mean as a linear regression on cycle number (e.g.,  $ k[A+D] $ for line  $ k $). This is also twice the average effect of substituting the unfavorable allele with the favorable one weighted by the allele-frequency change.

The final dominance term

\begin{equation}
D_{q}=-2\sum\left(\Delta p_{i}\right)^{2}d_{i}
\tag{18.39d}
\label{eq:18_39d}
\end{equation}

appears as a quadratic function of cycle number, with terms of the form  $ k^{2}D_{q} $ for cycle k and  $ k\dot{j}D_{q} $ for crosses between lines from cycles j and k. This is perhaps the most interesting dominance term, as it measures both inbreeding depression and the expected heterosis when two lines are crossed. As discussed below,  $ D_{q} $ has also been interpreted as the component of change in the mean due to the effects of drift, with  $ A + D $ representing the response from

625

selection (the expected deterministic change). Recall (LW Chapter 10) that the heterosis, H, between two lines

\begin{equation}
H=\mu_{1\times2}-\frac{\mu_{1}+\mu_{2}}{2}
\end{equation}

is the difference between the mean of the  $ F_{1} $ and the average of its parental lines. In the absence of epistasis

\begin{equation}
H=\sum(\delta p_{i})^{2}d_{i}
\tag{18.39e}
\label{eq:18_39e}
\end{equation}

where  $ \delta p_i $ is the between-line difference in allele frequency for the favorable allele at the  $ i $th locus and  $ d_i $ is the associated dominance term. Because the expected difference,  $ \delta p_i $, in allele frequency between lines separated by a single cycle of selection is  $ \Delta p_i $, comparing Equations 18.39d and 18.39e shows that  $ D_q $ is the negative of twice the heterosis generated by a single generation of selection, namely,  $ D_q = -2H $.

The composite parameters  $ (m, A, D, D_0, $ and  $ D_q) $ are typically estimated by OLS, with the vector of estimates given by  $ \left(\mathbf{X}^T\mathbf{X}\right)^{-1}\mathbf{X}^T\mathbf{y} $, with  $ \mathbf{y} $ being the vector of population means and the design matrix,  $ \mathbf{X} $, containing the appropriate coefficients for the parameters of interest (Example 18.9). The machinery introduced in LW Chapter 9 for standard line-cross analysis, such as hypothesis testing, testing goodness-of-fit, and estimation via GLS when the standard errors vary over means, can all be applied in a straightforward fashion. As with a standard line-cross analysis, one starts a GMA with the simple additive model and then tests whether adding additional model terms significantly improves the fit (LW Chapter 9). Typically one ends up with a reduced model, as some of the parameters turn out to be nonsignificant.

GMA can be extended to situations where some of the crosses are between lines from different base populations (Smith 1979b, 1983; Helms et al. 1989; Butruille et al. 2004; Melchinger and Flachenecker 2006). The same logic holds as before, with the frequency of the favorable allele in cycle $k$ from base population $j$ being $p_j(k) \simeq p_j + k \Delta p_j$, and with associated parameters $m_j, A_j, D_j, D_{0,j}$, and $D_{q,j}$. Additional terms are required to account for any initial heterosis (cycle-zero crosses between base populations $j$ and $i$) which, from Equation 18.39e, is a function of the square of the difference in their allele frequencies, $\delta_{i,j}^2 = (p_i - p_j)^2$. Finally, terms for any additional heterosis that has accrued during selection are required, and again these are functions of the squared allele-frequency differences. For a cycle-$k$ line from base population $i$ crossed to a cycle-$l$ line from base population $j$, this is

\begin{equation}
[\delta_{i,j}(k,\ell)]^{2}=[p_{j}(k)-p_{i}(\ell)]^{2}=[(p_{j}+k\Delta p_{j})-(p_{i}+\ell\Delta p_{i})]^{2}
\end{equation}

The motivation for this type of analysis is examined in our final volume, where we discuss the breeding of superior hybrids between lines. Basically, one tries to improve the performance of both selected lines, as well as that of their resulting hybrid. Generation means analysis can help partition any observed improvement into components of within-line and between-line (hybrid) improvement.

\begin{table}[h]
\centering
\begin{tabular}{|c|c|}
\hline
Cell 1 & Cell 2 \\
\hline
Cell 3 & Cell 4 \\
\hline
\end{tabular}
\caption{Table Placeholder}
\end{table}

626

\begin{table}[h]
\centering
\begin{tabular}{|c|c|}
\hline
Cell 1 & Cell 2 \\
\hline
Cell 3 & Cell 4 \\
\hline
\end{tabular}
\caption{Table Placeholder}
\end{table}

The resulting vector, y, of means and the design matrix, X, become

\begin{equation}
\mathbf{y}=\begin{pmatrix}{{{59.7}}} \\{{{63.4}}} \\{{{68.9}}} \\{{{42.3}}} \\{{{61.2}}} \\{{{79.6}}} \\{{{72.3}}} \\{{{40.3}}}\end{pmatrix},\qquad\mathbf{X}=\begin{pmatrix}{{{1}}}&{{{0}}}&{{{1}}}&{{{0}}} \\{{{1}}}&{{{4}}}&{{{1}}}&{{{4}}} \\{{{1}}}&{{{7}}}&{{{1}}}&{{{7}}} \\{{{1}}}&{{{7}}}&{{{0.5}}}&{{{3.5}}} \\{{{1}}}&{{{2}}}&{{{1}}}&{{{2}}} \\{{{1}}}&{{{3.5}}}&{{{1}}}&{{{3.5}}} \\{{{1}}}&{{{5.5}}}&{{{1}}}&{{{5.5}}} \\{{{1}}}&{{{5.5}}}&{{{0.5}}}&{{{2.75}}}\end{pmatrix}
\end{equation}

Solving for the vector of parameter estimates yields

\begin{equation}
\begin{aligned}\left(\begin{array}{c}\widehat{m}\\\widehat{A}\\\widehat{D}_{0}\\\widehat{D}\\\end{array}\right)=\left(\mathbf{X}^{T}\mathbf{X}\right)^{-1}\mathbf{X}^{T}\mathbf{y}=\left(\begin{array}{c}4.266\\1.072\\57.402\\0.523\\\end{array}\right)\end{aligned}
\end{equation}

Based on this subset of data, we note that there is considerable initial dominance ( $ D_0 = 57.402 $) influencing the initial value of the trait. Of the per-generation change of  $ A + D = 1.595 $ per generation, roughly two-thirds are from  $ A $ (1.072), and one-third is from  $ D $ (0.5223). Standard errors for these estimates follow using standard OLS machinery (LW Chapter 9).

Example 18.10. Here we derive the line and cross means when dominance is present. Because the focus is on the contribution from a specific locus, for ease of presentation we suppress the subscript (on a, d, and p) denoting the locus being followed. Assuming no epistasis, the mean is derived by summing the contributions over all loci. To allow for dominance, the genotypic values at our target locus become 0 : a + d : 2a. With p being the frequency of the favorable allele, the mean contribution from this locus is

\begin{equation}
0\cdot(1-p^{2})+(a+d)\cdot2p(1-p)+(2a)\cdot p^{2}=2ap+2dp(1-p)
\end{equation}

Notice that the contribution from $a$ is a linear function of the allele frequency (and hence a linear function of allele frequency change), while the contribution from $d$ has both linear ($p$) and quadratic ($p^2$) components. The contribution from this locus to the line mean in cycle $k$ follows by replacing $p$ with $p(k) \simeq p + k \Delta p$ in Equation 18.36b to yield

\begin{equation}
\begin{aligned}C_{k}&=\mu+2ap(k)+2dp(k)[1-p(k)]\\&=\mu+2a[p+k\Delta p]+2d[p+k\Delta p][1-(p-k\Delta p)]\\&=\left(\mu+2ap\right)+k\cdot2a\Delta p+2dp(1-p)+k\cdot2d\Delta p\cdot(1-2p)-k^{2}\cdot2d\left(\Delta p\right)^{2}\\&=m+k\cdot A+D_{0}+k\cdot D+k^{2}\cdot D_{q}\\ \end{aligned}
\end{equation}

Now suppose that this line is selfed. Homozygotes replicate themselves, while segregation of genotypes occurs when heterozygotes are selfed, with half of the offspring being  $ a + d $ heterozygotes, one-quarter of the offspring being 2a homozygotes, and the final one-quarter

627

being 0-valued homozygotes. The resulting mean value following a generation of selfing is the sum of these contributions,

\begin{equation}
p^{2}\left[2a\right]+2p(1-p)\left[(1/4)(2a)+(1/2)(a+d)+(1/4)\cdot0\right]+(1-p^{2})\cdot0=2ap+dp(1-p)
\end{equation}

Note that all terms involving d are exactly half of their random-mating values, yielding

\begin{equation}
C_{k}(s)=m+k\cdot A+\frac{1}{2}\left(D_{0}+k\cdot D+k^{2}\cdot D_{q}\right)
\end{equation}

Now consider the mean of the cross of lines from cycles k and j. As above, the chance of a 0-valued homozygote is  $ [1-p(k)][1-p(j)] $, whereas it is  $ p(k)[1-p(j)]+p(j)[1-p(k)] $ for an  $ a+d $ heterozygote, and  $ p(k)p(j) $ for a 2a homozygote, yielding a resulting mean of

\begin{equation}
\begin{aligned}C_{k\times j}&=\mu+(a+d)\cdot\left(p(k)\left[1-p(j)\right]+p(j)\left[1-p(k)\right]\right)+2a\cdot p(k)p(j)\\&=\mu+a\Biggl(p(k)+p(j)\Biggr)+d\cdot\Biggl(p(k)+p(j)-2p(k)p(j)\Biggr)\\=\left(\mu+2ap\right)&+(k+j)a\Delta p+d\Biggl(2p(1-p)+(k+j)\cdot\Delta p(1-2p)-k j\cdot2\left(\Delta p\right)^{2}\Biggr)\\&=m+\frac{k+j}{2}\cdot A+D_{0}+\frac{k+j}{2}\cdot D+k j\cdot D_{q}\end{aligned}
\end{equation}

If we self these individuals to form an  $ F_{2} $, by the same argument used above for selfing a line, each of the dominance terms is reduced by one-half, recovering Equation 18.38d.

Finally, when we form an  $ F_{2} $ by randomly mating the  $ F_{1} $ progeny from a line cross, the new allele frequency is

\begin{equation}
\frac{p(k)+p(j)}{2}\simeq p+\frac{k+j}{2}\Delta p
\end{equation}

yielding a mean of

\begin{equation}
\begin{align*}C_{k\times j}(rm)&=\mu+2a\left(p+\frac{k+j}{2}\Delta p\right)+2d\left(p+\frac{k+j}{2}\Delta p\right)\left(1-p-\frac{k+j}{2}\Delta p\right)\\=(\mu+2ap)+\frac{k+j}{2}2a\Delta p+2dp(1-p)+\frac{k+j}{2}2d\Delta p(1-2p)-\left(\frac{k+j}{2}\right)^{2}2d(\Delta p)^{2}\\&=m+\frac{k+j}{2}\cdot A+D_{0}+\frac{k+j}{2}\cdot D+\left(\frac{k+j}{2}\right)^{2}\cdot D_{q}\end{align*}
\end{equation}

Recalling Equation 18.38c, we can express this as

\begin{equation}
C_{k\times j}(r m)=C_{k\times j}-k j\cdot D_{q}+\left(\frac{k+j}{2}\right)^{2}\cdot D_{q}=C_{k\times j}+\frac{(k-j)^{2}}{4}\cdot D_{q}
\end{equation}

The additional term is due to the decay in heterosis following random mating, as the heterosis in the  $ F_{1} $ cross between lines from cycles j and k is

\begin{equation}
\begin{aligned}H&=\sum(p_{i}(k)-p_{i}(j))^{2}d_{i}=\sum(p+k\Delta p-p-j\Delta p)^{2}d_{i}\\&=(k-j)^{2}\sum(\Delta p)^{2}d_{i}=-(k-j)^{2}D_{q}/2\end{aligned}
\end{equation}

Recall (LW Chapter 9) that the amount, H, of heterosis from the  $ F_{1} $ is reduced to  $ H/2 $ in the (random-mating generated)  $ F_{2} $. Hence

\begin{equation}
C_{k\times j}(rm)-C_{k\times j}=-\frac{H}{2}=(k-j)^{2}\frac{D_{q}}{4}
\end{equation}

628

as found above.

\section{Accounting for Inbreeding Depression and Drift}
The deterministic assumption of a constant rate of allele-frequency change in the Hammond-Gardner model requires large population sizes. Because a typical selection experiment often involves small populations, genetic drift can also be important and thus needs to be taken into account. Smith (1979a, 1979b, 1983) was the first to incorporate drift and inbreeding depression in a GMA, while the most general treatment is from Melchinger and Flachenecker (2006). The latter authors' key idea was to separate allele frequency at cycle  $ k $ into deterministic and drift components.

\begin{equation}
p(k)=p(0)+k\Delta p+\delta p_{k}
\end{equation}

where the  $ \Delta p $ represents the expected large-population change (i.e., the Hammond-Gardner model), while  $ \delta p_k $ accounts for the additional change due to drift. Recall from Chapter 2 that under drift  $ E[\delta p_k] = 0 $, so that  $ E\left[(\delta p_k)^2\right] = \sigma^2(\delta p_k) $. Under pure drift, Equation 2.14a yields an allele-frequency variance after  $ k $ generations of

\begin{equation}
\sigma^{2}(\delta p_{k})=p(0)[1-p(0)]f_{k}
\tag{18.41a}
\label{eq:18_41a}
\end{equation}

where  $ f_k \simeq k/(2N_e) $ is the amount of inbreeding in cycle  $ k $. Melchinger and Flachenecker suggested that, with selection, Equation 18.41a can be approximated by using the average frequency over the  $ k $ cycles

\begin{equation}
E[(\delta p_{k})^{2}]\simeq\left(\frac{p(0)[1-p(0)]+p(k)[1-p(k)]}{2}\right)f_{k}
\tag{18.41b}
\label{eq:18_41b}
\end{equation}

Using the deterministic approximation  $ p(k) = p(0) + k\Delta p $, this reduces to

\begin{equation}
E[\left(\delta p_{k}\right)^{2}]\simeq p(0)[1-p(0)]\;f_{k}+\frac{k}{2}\cdot(1-2p)\Delta p\;f_{k}-\frac{k^{2}}{2}\left(\Delta p\right)^{2}\;f_{k}
\tag{18.41c}
\label{eq:18_41c}
\end{equation}

Notice the close similarity of these three terms to  $ D_{0} $,  $ D $, and  $ D_{q} $ (c.f. Equations 18.39a–18.39d). Using Equation 18.40 as the model for allele frequency and following the logic used in Example 18.10, the adjustment in the mean for inbreeding due to drift becomes

\begin{equation}
I_{k}=\left(D_{0}+\frac{k\cdot D}{2}-\frac{k^{2}\cdot D_{q}}{2}\right)f_{k}
\tag{18.42a}
\label{eq:18_42a}
\end{equation}

Assuming  $ f_k = k \Delta f $ (typically, with  $ \Delta f = 1/[2N_e] $), this simplifies to

\begin{equation}
I_{k}=k\left(D_{0}+\frac{k\cdot D}{2}-\frac{k^{2}\cdot D_{q}}{2}\right)\Delta f
\tag{18.42b}
\label{eq:18_42b}
\end{equation}

The resulting line means are

\begin{equation}
C_{k}=C_{k}(HG)-I_{k}\qquadand\qquad C_{k}(s)=C_{k}(s,HG)-I_{k}/2
\tag{18.43a}
\label{eq:18_43a}
\end{equation}

while the line-cross means become

\begin{equation}
C_{k\times j}=C_{k\times j}(HG)-I_{k}\quad for j>k
\tag{18.43b}
\label{eq:18_43b}
\end{equation}

where HG denotes the value under the Hammond-Gardner model (Equations 18.38a–18.38e).

629

Equations 18.42 and 18.43 define the Melchinger-Flachenecker model. The presence of drift  $ (f_k > 0) $ modifies the contributions from the three dominance terms, resulting in a reduction in mean (relative to the HG model) when both  $ D_0 $ and  $ D $ are positive (as  $ D_q $ is  $ \leq 0 $ when directional dominance is present). Smith (1979a, 1983) equated this reduction to inbreeding depression, which arises from drift causing some favorable alleles to decrease in frequency. On average, drift is equally likely to increase or decrease the change in frequency of an allele over its deterministic value,  $ \Delta p $. Thus, under drift, some alleles have a greater-than-deterministic increase in frequency, but this is countered by (on average) others having a less-than-expected increase. For additive genes, the expected value of these changes exactly cancel (as  $ E[\delta p_k] = 0 $). However, with dominance, this variance in allele-frequency change around  $ \Delta p $ does not cancel out, as the mean is a quadratic function of  $ d $.

While Melchinger and Flachenecker’s analysis is exact, the approximate results of the work by Smith (1979a, 1983) had a significant influence on the GMA literature, so a few comments on his approximations are in order. Smith’s (1979a) model was

\begin{equation}
C_{k}(s,S)=C_{k}(s,HG)+k\cdot I
\tag{18.44}
\label{eq:18_44}
\end{equation}

namely, it added a term,  $ kI $, to the Hammond-Gardner model expression for  $ C_k(s) $ (Equation 18.38b) to account for inbreeding, with  $ I = (D_0 + D) $. Note by comparison to the exact result (Equation 18.42b) that this is only correct when  $ D_0 \gg D $, and  $ D_q $ is negligible (as Smith assumed). Further, Smith neglected to include  $ I $ in his expressions for the line-cross means, instead using the HG expression for  $ C_{k \times j}(s) $, Equation 13.38d. Equation 18.43b shows this to be incorrect. Smith (1979a) ignored the  $ D_q $ term, as he felt that  $ \Delta p $ was expected to be small, and hence its square was likely to be negligible (or in any case, potentially difficult to estimate with precision).

In a later paper, Smith (1983) instead ignored the inbreeding correction,  $ I $, and suggested that for experiments with very small effective population size,  $ E\left[\left(\delta p_k\right)^2\right] \gg \left(\Delta p\right)^2 $, and thus  $ D_q $ was a surrogate measure for drift, as most of its value was due to the drift variance. Equations 18.42b and 18.43 show that both of these ideas, while correct in spirit, are approximations of the more general solution (Equation 18.42b). However, they can often yield a fairly good approximation. For example, Flachenecker et al. (2006) found that both the Smith (1983) and Melchinger-Flachenecker models gave rather similar parameter estimates. The critical insight offered by Smith's analysis was that by adjusting for the effects of drift, one could estimate the potential response if a larger population size was used, as the following two examples illustrate.

Example 18.11. Helms et al. (1989) obtained an estimate of  $ \widehat{D}_q = -0.024 $ (which was significant at the 1% level) in an analysis of a maize line selected for yield. Assuming small effective population size, how much was the population mean at cycle ten reduced due to drift-generated inbreeding depression? Recalling from Equation 18.38a that the contribution from  $ D_q $ in cycle  $ k $ is  $ D_q \cdot k^2 = -0.024 \cdot 10^2 = -2.4 $, then the observed mean (provided the assumptions of the Smith [1983] model hold) is expected to be reduced by 2.4 units (relative to a randomly mating population) by the effects of drift generating inbreeding depression.

\begin{table}[h]
\centering
\begin{tabular}{|c|c|}
\hline
Cell 1 & Cell 2 \\
\hline
Cell 3 & Cell 4 \\
\hline
\end{tabular}
\caption{Table Placeholder}
\end{table}

630

\begin{table}[h]
\centering
\begin{tabular}{|c|c|}
\hline
Cell 1 & Cell 2 \\
\hline
Cell 3 & Cell 4 \\
\hline
\end{tabular}
\caption{Table Placeholder}
\end{table}

Finally, a rough calculation suggests the conditions under which  $ \sigma^2(\delta p) \gg (\Delta p)^2 $, i.e., when the variance in allele-frequency change from drift exceeds the square of the deterministic change from selection. Equation 5.2b shows that  $ \Delta p \simeq sp(1 - p) $ for a very weakly selected additive locus, while Equation 5.21 yields  $ s \simeq \bar{\imath} a / \sigma_z $, where  $ \bar{\imath} $ is the selection intensity on the trait, yielding  $ \Delta p_i \simeq (\bar{\imath} a_i / \sigma_z) p_i (1 - p_i) $. Recalling that  $ \sigma^2(A_i) = 2a_i^2 p_i (1 - p_i) $ is the additive genetic variance associated with locus  $ i $

\begin{equation}
\begin{align*}(\Delta p_{i})^{2}&\simeq(\bar{\imath}a_{i}/\sigma_{z})^{2}p_{i}^{2}(1-p_{i})^{2}=\bar{\imath}^{2}p_{i}(1-p_{i})\left(\frac{a_{i}^{2}p_{i}(1-p_{i})}{\sigma_{z}^{2}}\right)\\&=\bar{\imath}^{2}p_{i}(1-p_{i})\left(\frac{\sigma^{2}(A_{i})/2}{\sigma_{z}^{2}}\right)=\frac{\bar{\imath}^{2}p_{i}(1-p_{i})}{2}h_{i}^{2}\end{align*}
\end{equation}

where  $ h_{i}^{2} $ is the heritability due to locus i. Assuming that n loci underlying the trait contribute roughly the same to the overall heritability,  $ h^{2} $, then

\begin{equation}
(\Delta p_{i})^{2}\simeq\frac{p_{i}(1-p_{i})}{2}\frac{\bar{\imath}^{2}h^{2}}{n}
\tag{18.45a}
\label{eq:18_45a}
\end{equation}

Drift dominates the expected quadratic change when  $ \sigma^{2}(\delta p) \gg (\Delta p)^{2} $, or

\begin{equation}
\frac{p(1-p)}{2N_{e}}\gg\frac{p(1-p)}{2}\frac{\bar{\tau}^{2}h^{2}}{n}
\end{equation}

which reduces to

\begin{equation}
n\gg N_{e}\left(\bar{\imath}^{2}h^{2}\right)
\tag{18.45b}
\label{eq:18_45b}
\end{equation}

For example, for a trait with  $ h^2 = 0.25 $ under truncation selection to save the largest 5%  $ (\bar{i} = 2.06; $ Example 14.1 $ ) $,  $ \bar{i}^2 h^2 \simeq 1 $, so drift dominates the quadratic term when the number of loci is much greater than the effective population size (assuming loci have roughly equal effects). Expressed in another way, the fractional contribution to the expected square in allele-frequency change from drift is

\begin{equation}
\frac{\sigma^{2}(\delta p)}{\sigma^{2}(\delta p)+(\Delta p)^{2}}\simeq\frac{1/N_{e}}{1/N_{e}+\bar{\imath}^{2}h^{2}/n}=\frac{1}{1+\bar{\imath}^{2}h^{2}N_{e}/n}
\tag{18.45c}
\label{eq:18_45c}
\end{equation}

or

\begin{equation}
\frac{\sigma^{2}(\delta p)}{\sigma^{2}(\delta p)+(\Delta p)^{2}}\simeq1-\bar{\imath}^{2}h^{2}N_{e}/n,\quad\mathrm{w h e n}\quad n\gg N e
\tag{18.45d}
\label{eq:18_45d}
\end{equation}

For  $ h^2 = 0.25 $,  $ \bar{i} = 2.06 $,  $ n = 100 $, and  $ N_e = 25 $, Equation 18.45c returns a value of 0.79, while the approximation given by Equation 18.45d yields 0.74. Hence, in this setting almost 80% of the expected squared allele-frequency change is due to drift.

\end{document}